% concatenated from mist_2026.tex, sections/abstract.tex, sections/introduction.tex, sections/related_work.tex, sections/preliminaries.tex, sections/framework.tex, sections/theory.tex, sections/experiments.tex, sections/conclusion.tex, sections/CLiterature.tex, sections/appendix.tex
\documentclass{article}

% if you need to pass options to natbib, use, e.g.:
%     \PassOptionsToPackage{numbers, compress}{natbib}
% before loading neurips_2026

% The authors should use one of these tracks.
% Before accepting by the NeurIPS conference, select one of the options below.
% 0. "default" for submission
\usepackage[preprint]{neurips_2026}
% the "default" option is equal to the "main" option, which is used for the Main Track with double-blind reviewing.
% 1. "main" option is used for the Main Track
%  \usepackage[main]{neurips_2026}
% 2. "position" option is used for the Position Paper Track
%  \usepackage[position]{neurips_2026}
% 3. "eandd" option is used for the Evaluations & Datasets Track
 % \usepackage[eandd]{neurips_2026}
 % if you want to opt-in for a de-anomymized submission (not recommended, but OK):
 %\usepackage[eandd, nonanonymous]{neurips_2026}
% 4. "creativeai" option is used for the Creative AI Track
%  \usepackage[creativeai]{neurips_2026}
% 5. "sglblindworkshop" option is used for the Workshop with single-blind reviewing
 % \usepackage[sglblindworkshop]{neurips_2026}
% 6. "dblblindworkshop" option is used for the Workshop with double-blind reviewing
%  \usepackage[dblblindworkshop]{neurips_2026}

% After being accepted, the authors should add "final" behind the track to compile a camera-ready version.
% 1. Main Track
 % \usepackage[main, final]{neurips_2026}
% 2. Position Paper Track
%  \usepackage[position, final]{neurips_2026}
% 3. Evaluations & Datasets Track
 % \usepackage[eandd, final]{neurips_2026}
% 4. Creative AI Track
%  \usepackage[creativeai, final]{neurips_2026}
% 5. Workshop with single-blind reviewing
%  \usepackage[sglblindworkshop, final]{neurips_2026}
% 6. Workshop with double-blind reviewing
%  \usepackage[dblblindworkshop, final]{neurips_2026}
% Note. For the workshop paper template, both \title{} and \workshoptitle{} are required, with the former indicating the paper title shown in the title and the latter indicating the workshop title displayed in the footnote.
% For workshops (5., 6.), the authors should add the name of the workshop, "\workshoptitle" command is used to set the workshop title.
% \workshoptitle{WORKSHOP TITLE}

% "preprint" option is used for arXiv or other preprint submissions
 % \usepackage[preprint]{neurips_2026}

% to avoid loading the natbib package, add option nonatbib:
%    \usepackage[nonatbib]{neurips_2026}

\usepackage[utf8]{inputenc} % allow utf-8 input
\usepackage[T1]{fontenc}    % use 8-bit T1 fonts
\usepackage{hyperref}       % hyperlinks
\usepackage{url}            % simple URL typesetting
\usepackage{booktabs}       % professional-quality tables
\usepackage{amsfonts, amsmath, amssymb, amsthm}       % blackboard math symbols
\usepackage{nicefrac}       % compact symbols for 1/2, etc.
\usepackage{microtype}      % microtypography
\usepackage{xcolor}         % colors
\usepackage{caption}
\usepackage{subcaption}
\usepackage{graphicx}

% Algorithm
\usepackage{algorithm}
\usepackage{algorithmic}

% -------------------------------------------------------
% Theorem environments
% -------------------------------------------------------
\newtheorem{theorem}{Theorem}[section]
\newtheorem{proposition}[theorem]{Proposition}
\newtheorem{corollary}[theorem]{Corollary}
\newtheorem{lemma}[theorem]{Lemma}
\newtheorem{definition}[theorem]{Definition}
\newtheorem{remark}[theorem]{Remark}

% -------------------------------------------------------
% Macros
% -------------------------------------------------------
\newcommand{\mist}{\textsc{Mist}}
\newcommand{\vfdt}{\textsc{Vfdt}}
\newcommand{\efdt}{\textsc{Efdt}}
\newcommand{\slda}{\textsc{Slda}}
\newcommand{\sqda}{\textsc{Sqda}}
\newcommand{\arf}{\textsc{Arf}}
\newcommand{\Prob}{\mathbb{P}}


% Note. For the workshop paper template, both \title{} and \workshoptitle{} are required, with the former indicating the paper title shown in the title and the latter indicating the workshop title displayed in the footnote. 
\title{
    \mist{}: Reliable Streaming Decision Trees for Online Class-Incremental Learning via McDiarmid Bound
}


% The \author macro works with any number of authors. There are two commands
% used to separate the names and addresses of multiple authors: \And and \AND.
%
% Using \And between authors leaves it to LaTeX to determine where to break the
% lines. Using \AND forces a line break at that point. So, if LaTeX puts 3 of 4
% authors names on the first line, and the last on the second line, try using
% \AND instead of \And before the third author name.


\author{%
  Phu-Hoa Pham$^{\star}$\thanks{$^\star$\,Equal contribution.\quad
    $^\dagger$\,Corresponding author:
    \texttt{long.tran-thanh@warwick.ac.uk}}\\
  Faculty of Information and Technology \\
  University of Science\\
  Vietnam National University \\
  Ho Chi Minh City, Vietnam \\
  \texttt{23122030@student.hcmus.edu.vn} \\
  \And
  Chi-Nguyen Tran$^{\star}$ \\
  Faculty of Information and Technology \\
  University of Science\\
  Vietnam National University \\
  Ho Chi Minh City, Vietnam \\
  \texttt{23122044@student.hcmus.edu.vn} \\
  \And
  Phu-Quy Nguyen-Lam$^{}$ \\
  Faculty of Information and Technology \\
  University of Science\\
  Vietnam National University \\
  Ho Chi Minh City, Vietnam \\
  \texttt{23122048@student.hcmus.edu.vn}
  \And
  Dao Sy Duy Minh$^{}$ \\
  Faculty of Information and Technology \\
  University of Science\\
  Vietnam National University \\
  Ho Chi Minh City, Vietnam \\
  \texttt{23122041@student.hcmus.edu.vn}
  \And
  Huynh Trung Kiet$^{}$ \\
  Faculty of Information and Technology \\
  University of Science\\
  Vietnam National University \\
  Ho Chi Minh City, Vietnam \\
  \texttt{23122039@student.hcmus.edu.vn} \\
  \And
  Long Tran-Thanh$^{\dagger}$ \\
  Department of Computer Science \\
  University of Warwick \\
  Coventry, United Kingdom \\
  \texttt{long.tran-thanh@warwick.ac.uk}
}

\linespread{0.98}

\raggedbottom

\begin{document}
\maketitle
 
% -------------------------------------------------------
% ABSTRACT
% Style notes:
%   - No math formulas
%   - Follows: context -> gap -> what we do -> key result
%   - No bullet points, no overclaiming
% -------------------------------------------------------
\begin{abstract}
Streaming decision trees are natural candidates for open-world continual learning, as they perform local updates, enjoy bounded memory, and static decision boundaries. Despite these, they still fail in online class-incremental learning due to two coupled miscalibrations: (i) their split criterion grows unreliable as the class count $K$ expands, and (ii) the absence of knowledge transfer at split time. Both failures share a common root: the range of Information Gain intrinsically scales with $\log_2 K$. Consequently, any Hoeffding-style confidence radius derived from it must inevitably grow with the class count, making a $K$-independent split criterion structurally impossible, taking away the potential benefits of applying streaming decision trees to continual learning. To fix this issue, we present \mist{} (McDiarmid Incremental Streaming Tree), which resolves both failures through three integrated components: (i) a tight, $K$-independent McDiarmid confidence radius for Gini splitting that acts as a structural regulariser; (ii) a Bayesian inheritance protocol that projects parent statistics to child nodes via truncated-Gaussian moments, with variance reduction guarantees strongest precisely when splitting is most conservative; and (iii) per-leaf KLL quantile sketches that support both continuous threshold evaluation and geometry-adaptive leaf prediction from a single data structure. On standard and stress-test tabular streams, \mist{} is competitive with global parametric methods on near-Gaussian benchmarks and uniquely robust on non-Gaussian geometry where SOTA benchmarks collapse.
\end{abstract}
% -------------------------------------------------------
% INTRODUCTION
% Structure: broad context -> narrow to streaming trees ->
%   root cause (K-dependence) -> McDiarmid answer ->
%   cold start emerges -> "Against this background" -> contributions
% Style: - dense prose, "However" as pivot,
%   bold paragraph headers, numbered contributions at end
% -------------------------------------------------------
\section{Introduction}
Continual learning systems that operate in an open world must absorb new
classes as they arrive, without access to historical data and without
degrading what was previously learned. In tabular settings, where large
pretrained representations are unavailable, this challenge must be resolved
directly from the raw feature stream under strict memory constraints.

Streaming decision trees are a natural candidate for this regime. Their updates 
are local, as only the leaf reached by an arriving sample is modified. As such, their memory 
can remain bounded by design, and they share no global representation across classes. 
For example, in the foundational VFDT framework, which laid down the basis for streaming decision trees, a decision 
boundary, once committed via the Hoeffding bound, is not revised. Instead, 
splits are permanent, making catastrophic forgetting structurally 
avoidable in principle. This is an advantage against global parametric methods, as none of them can 
provide it by construction. This structural potential, however, 
has never been fully realised. In practice, classical streaming trees collapse in the online 
class-incremental learning (Online CIL) setting. The reason is more fundamental than it may 
first appear, and tracing it to its root reveals two coupled failures that must be 
resolved simultaneously.

\noindent
\textbf{The root cause: Premature split.}
Existing streaming trees typically 
trigger splits when the empirical gain gap
exceeds a Hoeffding confidence radius derived from Information Gain~\cite{DomingosHulten2000,ManapragadaWebbSalehi2018,korycki2021,BifetGavalda2009HAT,chaouki2024}. This
radius grows as $\log_2 K / \sqrt{n}$, which comes from the structural implication of
Information Gain's range scaling with $\log_2 K$, and becomes unreliable
as $K$ grows without bound, either delaying necessary splits or triggering
prematurely, making streaming decision trees not suitable for online CIL settings. 

\noindent
\textbf{A potential solution and its cold-start split issue.}
A natural way to fix this premature split issue is to switch to Gini
impurity, whose range is bounded in $[0,1)$ regardless of $K$, making a
class-count-independent radius achievable in principle. However, it is not far from trivial how to identify a \emph{optimal} decision radius
for the Gini split gain.
%
To answer this question, in this paper we derive a tight bounded-difference constant for Gini split gain under
single-sample replacement and prove it is unimprovable for arbitrary $K$,
yielding a tight class-count-independent confidence radius. 
%This is a qualitative change, not merely a tighter constant. 

This calibration is conservative
by design, acting as a structural regulariser. It does, however, expose
a second point of failure: When a leaf splits, its children inherit nothing from
their parent and must rebuild predictive distributions from scratch,
producing unreliable predictions precisely when the tree adapts most
aggressively; further splits from sparse children compound the
instability.

\noindent
\textbf{The coupled failures: premature and cold-start splits.}
The two failures reinforce each other, as premature splits produce sparse
children, whose cold-start inflates impurity estimates, triggering
further premature splits. To fix this, we break this cycle by resolving both
simultaneously: Conservative splitting ensures children are born from
data-rich parents, maximising the statistical value of what is inherited.
To do so, we use Bayesian inheritance and translate those statistics into tight posterior
distributions from the first sample onward. Crucially, variance reduction
from inheritance scales with the parent sample count at split time, so
the more conservative the criterion, the stronger the inheritance
guarantee.

\noindent
\textbf{Sketching for both splitting and prediction.}
The two components above handle \emph{when} to split and \emph{how} to
initialise after splitting. 
However, they need to rely on past raw data, which is typically not feasible in the online CIL setting.
To fix this, we propose a new, sketching based approach to handle continuous threshold evaluation
without retaining raw data. In particular, per-leaf, label-aware quantile sketches generate
candidate thresholds from empirical class-conditional distributions, and the
same structure powers geometry-adaptive leaf prediction---parametric Gaussian
on near-Gaussian data, distribution-free otherwise. Both configurations are
formalised in Section~\ref{sec:framework}.

\noindent
\textbf{Contributions.} In this paper, we present \mist{} (McDiarmid Incremental Streaming Tree) to implement the three abovementioned novel ideas as follows:
\begin{itemize}
  \item \textbf{Tight McDiarmid calibration for Gini splitting.}
  We derive the tight bounded-difference constant for multi-class Gini split
  gain and prove that it is sharp for arbitrary $K$, yielding a
  class-count-independent confidence radius (Corollary~\ref{cor:radius}).
  This calibration eliminates the $K$-scaling instability of Hoeffding-based
  criteria, a qualitative improvement, not merely a tighter constant, and
  acts as a structural regulariser protecting the inheritance mechanism.
  \item \textbf{Bayesian post-split stabilisation.}
  We introduce parent-to-child projection via exact truncated-Gaussian
  conditional moments for the split feature and Naive Bayes mass rescaling for
  remaining features. Variance reduction guarantees cover both
  class-probability estimates 
  and feature-parameter estimates
  (Proposition~\ref{prop:variance}) 
  with both
  posteriors tightening in proportion to the parent sample count at split time.
  \item \textbf{Unified label-aware sketching for structure and prediction.}
  Per-leaf KLL quantile sketches support continuous split threshold evaluation
  without raw-data retention; the same sketches power an optional
  non-parametric leaf predictor under a Naive Bayes factorisation at no
  additional storage cost. Theorem~\ref{thm:composed} provides a composed
  error bound jointly controlling statistical variance and sketch rank error.
  \item \textbf{Regime-aware empirical characterisation.}
  We identify and explain conditions under which the parametric and
  non-parametric leaf predictors are respectively preferable, supported by
  prequential accuracy curves, ablation studies, and stress-test experiments
  on synthetic streams with non-Gaussian class geometry.
\end{itemize}
\section{Related Work}
In this section due to the space limits we only cover the most relevant work to our paper, and defer the more detailed summary of the online continual learning literature to Appendix~\ref{app:related_work}. 

\noindent
\textbf{Streaming decision trees in class-incremental settings.}
\vfdt{}~\cite{DomingosHulten2000} established Hoeffding-style confidence tests for incremental splitting, with EFDT~\cite{ManapragadaWebbSalehi2018} adding continuous re-evaluation and HAT~\cite{BifetGavalda2009HAT} incorporating adaptive drift detection via ADWIN~\cite{BifetGavalda2007ADWIN}. These methods treat distributional shifts as concept drift to be discarded, directly conflicting with Online CIL where class arrivals are permanent structural events; disabling their drift detectors recovers substantial accuracy on our benchmarks, confirming the failure is protocol misalignment rather than algorithmic. \vfdt{} zero-initialises child leaves; \textsc{VFDTc}~\citep{Gama2003} introduced Naive Bayes leaves but retained zero-initialisation at split time. Korycki and Krawczyk~\cite{korycki2021} motivate knowledge transfer at structural boundaries under replay; our work addresses the analogous problem under the strict exemplar-free constraint with formal variance reduction guarantees. Thompson-sampling trees~\cite{chaouki2024} are restricted to categorical attributes and a closed-world class set, making direct comparison inappropriate.

\noindent
\textbf{Concentration bounds for Gini-based splitting.}
Rutkowski et al.~\cite{rutkowski2013} established that McDiarmid's inequality applies to Gini split gain with constant $c_i = 8/n$---the first $K$-independent confidence radius. De Rosa and Cesa-Bianchi~\cite{derosa2015splitting} tightened this to $4/n$ for $K=2$ only. Our work closes the remaining gap: $4/n$ is tight for arbitrary $K$, so the radius does not degrade as new classes arrive. C-Tree~\cite{derosa2015splitting} derives refined large-deviation intervals using problem-dependent parameters, but its analysis is restricted to $K=2$; an empirical comparison is omitted as extending its implementation to $K>2$ is non-trivial.

\noindent
\textbf{Global methods for exemplar-free tabular CIL.}
\slda{}~\cite{hayes2020lifelong} uses a shared-covariance Gaussian discriminant, \sqda{}~\cite{hayes2020lifelong} relaxes
this to per-class Gaussians, NCM~\cite{mensink2013} tracks class centroids, and
RLS~\cite{zhuang2022acil} maintains a shared linear classifier via
recursive least squares. These are competitive on near-Gaussian
benchmarks but share a structural vulnerability: all classes are coupled
through one global representation, so updates for new classes gradually
displace boundaries learned for earlier ones. Streaming decision trees
avoid this coupling by construction, and on non-Gaussian geometry where
global parametric methods collapse, a well-calibrated tree with
geometry-adaptive prediction remains competitive.

\noindent
\textbf{Non-parametric and sketch-based prediction in streaming settings.}
KDE Naive Bayes~\cite{john1995,Moghaddam2002} improves accuracy under non-Gaussian
class-conditionals, but raw-sample KDE is infeasible in an exemplar-free
stream. Prior work~\cite{benhaim2010} uses additive-error quantile sketches for split threshold
evaluation without connecting the sketch to the leaf
predictor. Our work unifies both roles: the same per-leaf KLL sketches~\cite{Karnin2016} that
support continuous split evaluation also power a distribution-free leaf
predictor, yielding a streaming analog of KDE Naive Bayes with formal error
guarantees. KLL’s efficient update mechanism maintains a strict memory bound, and its high-probability rank guarantees feed directly into Theorem~\ref{thm:composed}, with $O(k)$ memory and
$O(\log k)$ amortised update cost.
\section{Preliminaries and Problem Setting}
\label{sec:prelim}
\subsection{Strict Online CIL}
We consider learning from an infinite stream $\mathcal{S} = \{(x_t, y_t)\}_{t=1}^\infty$
where the label space $\mathcal{Y}_t$ expands monotonically and
$K = \lim_{t\to\infty}|\mathcal{Y}_t|$ is unbounded and unknown a priori.
An algorithm operates in the \emph{strict Online CIL} regime if it satisfies:
(i)~single-pass processing; (ii)~exemplar-free memory; (iii)~anytime classification.
New-class arrivals are permanent structural events, not distributional shifts
to be discarded. Each split decision is evaluated locally on the $n$ samples
accumulated at a leaf; class-conditional distributions $p(x \mid y=c)$ are
stationary by protocol, and samples at a leaf are treated as an independent
batch~\cite{DomingosHulten2000}, satisfying McDiarmid's independence requirement.
\enlargethispage{2\baselineskip}
\subsection{Gini Split Gain and Bounded Differences}
Let $S$ be $n$ samples at a leaf with empirical class proportions $p_c = n_c/n$.
The Gini impurity and split gain for a binary partition $(S_L, S_R)$ are:
\begin{equation}
  G(S) = 1 - \sum_{c=1}^K p_c^2, \qquad
  F(S) = G(S) - \tfrac{n_L}{n}G(S_L) - \tfrac{n_R}{n}G(S_R).
  \label{eq:gain}
\end{equation}
A function $f : (\mathcal{X} \times \mathcal{Y})^n \to \mathbb{R}$ satisfies the \emph{bounded difference
property} with constants $c_1,\ldots,c_n$ if replacing any single input changes
$f$ by at most $c_i$. McDiarmid's inequality then gives:
\begin{equation}
\Prob\!\left[\,|f(S) - \mathbb{E}[f]| \geq \varepsilon\right]
\leq 2\exp\!\left(-\frac{2\varepsilon^2}{\sum_i c_i^2}\right).
\label{eq:mcdiarmid}
\end{equation}
The central theoretical question of this paper is: what is the tight value
of $c_i$ for the Gini split gain $F$? The answer, derived in
Section~\ref{sec:theory}, determines the $K$-independent confidence radius
used by \mist{}.
% -------------------------------------------------------
% THE MIST FRAMEWORK
% Restructured order: split criterion -> inheritance -> sketching
% Rationale: McDiarmid-first narrative. Conservative split exposes cold start,
%   which motivates inheritance. Sketching then unifies both.
% -------------------------------------------------------
\section{The \mist{} Framework}
\label{sec:framework}
Our stream decision tree solution, \mist{}, is built around three tightly integrated components: (i) a split criterion determining when to partition a leaf; (ii) an inheritance mechanism initialising children with statistically meaningful priors; and (iii) a sketching component providing both continuous threshold evaluation and optional geometry-adaptive prediction from a single data structure.
\subsection{Sensitivity-Calibrated Split Criterion}
\label{sec:split}
\mist{} triggers a split when the empirical Gini gain gap $\widehat{\Delta F}$
between the best and second-best candidate feature exceeds a confidence
radius. The radius is derived from the tight, class-count-independent McDiarmid calibration in
Section~\ref{sec:theory} and is class-count-independent.

\begin{corollary}[Operational McDiarmid Radius]
\label{cor:radius}
Given failure probability $\delta$, $d$ candidate features, and $m$
candidate thresholds per feature, \mist{} triggers a split if:

\begin{equation}
  \widehat{\Delta F} > \varepsilon, \qquad
  \varepsilon = \sqrt{\frac{32\ln(2dm/\delta)}{n}}.
  \label{eq:radius}
\end{equation}
\end{corollary}

\begin{remark}
The factor of 32 arises from the GapTest structure: concluding
$\mathbb{E}[F_1] > \mathbb{E}[F_2]$ requires controlling both tails
simultaneously, so $\varepsilon = 2\eta$ where $\eta =
\sqrt{8\ln(2dm/\delta)/n}$ is the per-tail McDiarmid radius, giving
$\varepsilon = \sqrt{32\ln(2dm/\delta)/n}$. The union bound over $d$
features and $m$ thresholds enters only additively inside the logarithm
as $\ln(2dm/\delta)$, not multiplicatively into the constant 32.
% Crucially, $\varepsilon$ does not depend on $K$, whereas Hoeffding-based
% criteria scale as $\log_2 K$.
\end{remark}

Crucially, $\varepsilon$ does not depend on the global class count $K$,
whereas Hoeffding-based criteria scale as $\log_2 K$; the
threshold count $m$ from Section~\ref{sec:sketch} contributes only a
mild $O(\sqrt{\ln m})$ overhead.

\noindent
\textbf{Role as structural regulariser.}
The conservatism of this criterion is a feature:
% , not a limitation:
replacing the McDiarmid GapTest with a more aggressive criterion
increases tree size substantially at the cost of only marginal accuracy gains
(Table~\ref{tab:ablation-core}), consistent with the coupled failure
mode analysis (Figure~\ref{fig:tree-size}).
% The conservatism of this criterion is a feature:
% % , not a limitation:
% replacing the McDiarmid GapTest with a more aggressive criterion
% increases tree size substantially without improving accuracy
% (Table~\ref{tab:ablation-core}), consistent with the coupled failure
% mode analysis (Figure~\ref{fig:tree-size}).

\subsection{Bayesian Knowledge Inheritance}
\label{sec:inheritance}
\noindent
\textbf{Motivation.}
A leaf that has accumulated sufficient evidence to split has rich
class-conditional statistics that should not be discarded. Standard streaming
trees initialise child leaves without prior information, forcing them to reconstruct predictive
distributions from scratch. In a strict exemplar-free setting, this
cold-start phase can be severe: The children must accumulate enough new
samples to form reliable predictions before they are useful, and any further
splits triggered from sparse children compound the problem. The conservative
split criterion of Section~\ref{sec:split} mitigates this by ensuring that
splits occur only from data-rich parents. We use Bayesian knowledge inheritance
to complete the picture by ensuring that the resulting children begin with
statistically grounded priors.

\noindent
\textbf{Truncated-Gaussian projection at split time.}
\mist{} uses Gaussian Naive Bayes (GNB) leaves, parameterised by
per-class, per-feature Welford accumulators $(\mu_j^c, \sigma_j^c)$ and
Dirichlet class-probability pseudo-counts. When a parent splits on
feature $j$ at threshold $v$, let $\zeta = (v - \mu_j^c)/\sigma_j^c$.
Exact truncated-Gaussian moment projection initialises the split-feature
statistics as:
\begin{align}
\label{eq:mu_left_right}
  \mu_j^{\mathrm{left},c}  = \mu_j^c
    - \sigma_j^c\frac{\phi(\zeta)}{\Phi(\zeta)}, \qquad
  \mu_j^{\mathrm{right},c} = \mu_j^c
    + \sigma_j^c\frac{\phi(\zeta)}{1 - \Phi(\zeta)},
\end{align}
with corresponding variance updates given in
Appendix~\ref{app:truncgauss}. For every non-split feature $k \neq j$,
under the Naive Bayes conditional independence assumption
$p(X_k \mid c, X_j \leq v) = p(X_k \mid c)$, parameters are inherited
unchanged. 
Class mass is reapportioned as
$n_c^{\mathrm{left}} \leftarrow \alpha\,\Phi(\zeta)\,n_c^{\mathrm{parent}}$
and
$n_c^{\mathrm{right}} \leftarrow \alpha\,(1{-}\Phi(\zeta))\,n_c^{\mathrm{parent}}$,
with Dirichlet pseudo-counts reweighted identically by factor $\alpha$.
% Class mass is reapportioned as
% $n_c^{\mathrm{left}} \leftarrow \Phi(\zeta)\,n_c^{\mathrm{parent}}$ and
% $n_c^{\mathrm{right}} \leftarrow (1-\Phi(\zeta))\,n_c^{\mathrm{parent}}$,
% with Dirichlet pseudo-counts reweighted identically.
 
For categorical features encoded as one-hot indicators, the split is
binary and no moment projection is required: each child $s \in \{0,1\}$
inherits class mass $n_c^s \leftarrow \alpha\, n_c(s)$, where
$n_c(s) = \#\{i : y_i = c,\, x_i^{(\mathrm{attr})} = s\}$
is the empirical co-occurrence count at the parent leaf,
$\alpha \in (0,1]$ is the inheritance discount
(Section~\ref{sec:theory-var}), and Laplace smoothing at prediction
time ensures non-zero posteriors for all classes.

The inherited prior decays at rate $O(1/n_{\mathrm{child}})$ as new
samples accumulate, allowing each child to specialise to its local
geometry while retaining the benefit of the parent's experience during
the critical early phase. 
% This projection is shared by both leaf
% predictor configurations introduced in Section~\ref{sec:sketch}: even
% in the non-parametric configuration, Gaussian moments are used solely
% to initialise child priors and apportion class mass immediately after a
% split, keeping post-split stabilisation and leaf prediction decoupled.
This projection applies to both configurations (Section~\ref{sec:sketch}):
Gaussian moments initialise class-probability priors and apportion class
mass after every split. \mist{}-G inherits the full Gaussian leaf;
\mist{}-K's KLL sketches start empty (Algorithm~\ref{alg:split}), so
feature likelihoods require new samples, though Dirichlet priors are
immediately available.

\noindent
\textbf{Joint design with the split criterion.}
The variance reduction from inheritance scales inversely with the parent
sample count at split time
(Proposition~\ref{prop:variance}): a
split triggered too early yields a weaker prior and larger cold-start
variance for the child. The McDiarmid GapTest ensures splits occur only
from sufficiently data-rich parents, maximising the statistical value
of the inherited statistics. The two mechanisms reinforce each other:
the more conservative the split criterion, the richer the parent
statistics at split time, and therefore the stronger the variance
reduction guarantee. 

\subsection{Label-Aware Sketching and the Unified Dual Role of KLL}
\label{sec:sketch}
\looseness=-1
\mist{} maintains per-leaf, per-class KLL quantile sketches of capacity
$k$. A sketch for class $c$ at leaf $\ell$ is created only when samples
of class $c$ are routed there, giving per-leaf memory
$O(d \cdot K_{\mathrm{local}}^{(\ell)} \cdot k)$. As splits refine
the partition and leaves become purer, $K_{\mathrm{local}}^{(\ell)}$
decreases toward 1, providing a natural self-clearing mechanism that
bounds memory in practice; a deterministic worst-case bound is
$O(dkL_tK_t)$, which is empirically loose because most leaves satisfy
$K_{\mathrm{local}}^{(\ell)} \ll K_t$ (Appendix~\ref{app:memory-scaling}).

\noindent
\textbf{Split threshold evaluation.}
Candidate thresholds are midpoints between adjacent class-conditional
medians from active sketches, yielding at most $m \leq K_{\mathrm{local}}-1$
candidates per feature, exactly the $m$ in Eq.~\eqref{eq:radius}, so the
union bound covers only candidates actually considered. For categorical
attributes encoded as one-hot indicators, $m \leq 1$ and
Eq.~\eqref{eq:radius} remains valid without modification.

% For categorical attributes, we apply one-hot encoding before streaming. Each
% one-hot feature is binary, so candidate evaluation is equivalent to a single
% threshold test (effectively $m \leq 1$), and the confidence radius in
% Eq.~\eqref{eq:radius} remains unchanged.
\noindent
\textbf{Leaf prediction: \mist{}-G and \mist{}-K.}
The same per-class sketches that generate split thresholds power an
optional non-parametric leaf predictor. In the default configuration
(\mist{}-G), leaf likelihoods are computed from the Gaussian parameters
maintained by the Welford accumulators; sketches contribute only to
split threshold evaluation. In the non-parametric configuration
(\mist{}-K), leaf likelihoods are instead estimated from sketch-derived
per-feature CDFs under a Naive Bayes factorisation:
\begin{equation}
  \hat{p}(x_j \mid c)
  = \frac{\widehat{F}_{j,c}(x_j + h_j)
          - \widehat{F}_{j,c}(x_j - h_j) + \epsilon_s}
         {2h_j + \epsilon_s},
  \label{eq:mistk-likelihood}
\end{equation}
where $\widehat{F}_{j,c}$ is the sketch-derived CDF, $h_j$ is the
inter-quantile spacing at $x_j$ (local bandwidth), and $\epsilon_s > 0$
is a dimensionless smoothing regulariser that prevents division by zero
when $h_j \to 0$; its contribution is negligible at interior points
where $h_j \gg \epsilon_s$.
% where $\widehat{F}_{j,c}$ is the sketch-derived CDF, $h_j$ is the
% inter-quantile spacing at $x_j$ (local bandwidth), and $\epsilon_s > 0$
% is a smoothing constant. This is a streaming analogue of KDE-based
% Naive Bayes~\cite{john1995,Moghaddam2002}: it approximates the
% class-conditional density via a finite-difference CDF estimate rather
% than a kernel sum over stored samples, and therefore requires no raw-data
% retention. 
% For categorical attributes, \mist{}-K uses a Bernoulli
% likelihood estimated from Dirichlet pseudo-counts, since binary support
% renders quantile-based density estimation degenerate.

\noindent
\textbf{Regime-dependent behaviour.}
\mist{}-G is preferable when class-conditional distributions are
approximately unimodal Gaussian, a condition empirically indicated by
strong \slda{}/\sqda{} performance on the same benchmark. \mist{}-K
is preferable under strongly non-Gaussian geometry (multimodal, angular,
or heavily skewed), where the Gaussian leaf assumption introduces
systematic bias that the sketch-based marginal estimator avoids.
The dual role of the KLL sketch is not an engineering convenience: 
the same data structure that makes continuous split evaluation 
memory-feasible also captures the distributional information needed for 
geometry-adaptive prediction.
% -------------------------------------------------------
% Algorithms
% -------------------------------------------------------
\begin{algorithm}[t]
\caption{\mist{}: Update and Predict}
\label{alg:mist}
\begin{algorithmic}[1]
\REQUIRE Sample $(x_t, y_t)$, tree $\mathcal{T}$, discount $\alpha$,
         failure probability $\delta$, grace period $g$
\STATE Route $x_t$ to leaf $\ell$
\STATE Update Welford accumulators $(\mu_j^c, \sigma_j^c)$ and Dirichlet
       pseudo-counts at $\ell$ with $(x_t, y_t)$
\STATE Update per-class KLL sketch at $\ell$ with $x_{t,j}$ for each
       feature $j$
\IF{$n_\ell \geq g$ \AND \textsc{GapTest}$(\ell, \delta)$}
  \STATE $j^*, v^* \leftarrow$ best feature and threshold from sketch medians
  \STATE \textsc{Split}$(\ell, j^*, v^*, \alpha)$
         \COMMENT{Algorithm~\ref{alg:split}}
\ENDIF
\STATE \textbf{return} $\arg\max_c \hat{p}(c \mid x_t, \ell)$
\STATE
\STATE \textbf{function} \textsc{GapTest}$(\ell, \delta)$
\STATE \quad $m \leftarrow K_{\mathrm{local}}^{(\ell)} - 1$
       \COMMENT{at most $K_{\mathrm{local}}^{(\ell)}-1$ candidate 
       thresholds per feature; see Section~\ref{sec:sketch}}
\STATE \quad $\varepsilon \leftarrow \sqrt{32\ln(2dm/\delta)/n_\ell}$
\STATE \quad \textbf{return} $\widehat{F}_{(1)} - \widehat{F}_{(2)} > \varepsilon$
\end{algorithmic}
\end{algorithm}

% -------------------------------------------------------
% THEORETICAL ANALYSIS
% -------------------------------------------------------
\section{Theoretical Analysis}
\label{sec:theory}
The central result of this section is Theorem~\ref{thm:tight}: A tight,
$K$-independent bounded-difference constant for Gini split gain.
This is what makes Corollary~\ref{cor:radius} valid in an open world
and what Hoeffding-based criteria fundamentally cannot provide.
Theorem~\ref{thm:composed} extends this to the sketched setting.
Propositions~\ref{prop:variance}
quantify variance reduction from Bayesian inheritance.
\subsection{Tightness of the Gini Sensitivity}
\label{sec:theory-tight}
\begin{theorem}[Tightness of the Gini Sensitivity]
\label{thm:tight}
Let $S$ be a sequence of $n \geq 2$ samples and $F(S)$ the Gini split
gain defined in Eq.~\eqref{eq:gain}. For any single-sample replacement
$z_i \to z_i'$, the bounded-difference constant satisfies $c_i \leq 4/n$.
Moreover, this rate is tight: for any $\gamma > 0$, there exists a
configuration $S$ and a replacement such that
$|F(S) - F(S^{(i)})| > 4/n - \gamma$.
\end{theorem}

% \begin{proof}[Proof sketch]
% Case~A (same-child routing): replacing $z_i$ with $z_i'$ in child $L$
% leaves $n_L$, $n_R$, and $G(S_R)$ unchanged, so only $G(S)$ and
% $G(S_L)$ vary; each changes by at most $2/n$ and $2/n_L$, giving
% $|\Delta F| \leq 2/n + (n_L/n)(2/n_L) = 4/n$.
% Case~B (cross-routing): the replacement reroutes $z_i$ between children,
% changing all five quantities $n_L, n_R, G(S), G(S_L), G(S_R)$
% simultaneously. Expanding $n\Delta F$ into three signed terms and
% bounding each as a quadratic in the routing counts yields
% $n|\Delta F| < 2 < 4$ (full derivation in Appendix~\ref{app:case-b}).
% Tightness: take binary labels, $m = \lfloor\sqrt{n/2}\rfloor$
% left-child samples (all class~1) and $n{-}m$ right-child samples
% (all class~0); flip one left sample to class~0.
% The resulting deviation equals $H(m,n)/n$ where
% $H(m,n) = 2(n-m)(2m-1)/(nm)$.
% Since $m = o(n)$ and $m \to \infty$, both $(n-m)/n \to 1$ and
% $(2m-1)/m \to 2$, giving $H(m,n) \to 4$, so
% $|\Delta F| \to 4/n$ as $n \to \infty$.
% \end{proof}
\begin{proof}[Proof sketch]
Case~A (same-child routing): replacing $z_i$ with $z_i'$ in child $L$
leaves $n_L$, $n_R$, and $G(S_R)$ unchanged, so only $G(S)$ and
$G(S_L)$ vary; each changes by at most $2/n$ and $2/n_L$, giving
$|\Delta F| \leq 2/n + (n_L/n)(2/n_L) = 4/n$.
Case~B (cross-routing): direct expansion over disjoint class proportions
gives $n|\Delta F| \leq 2 < 4$ (Appendix~\ref{app:case-b}).
Tightness: binary labels with $\lfloor\sqrt{n/2}\rfloor$ left-child
samples yield $|\Delta F| \to 4/n$ as $n \to \infty$.
\end{proof}

\begin{remark}
\label{rem:k-independence}
Rutkowski et al.~\cite{rutkowski2013} showed $c_i = 8/n$ for arbitrary $K$;
De Rosa and Cesa-Bianchi~\cite{derosa2015splitting} tightened this to $4/n$
for $K=2$ only, without establishing tightness.
Theorem~\ref{thm:tight} closes the gap: $4/n$ is tight for all $K$,
so the operational radius of Corollary~\ref{cor:radius} does not degrade
as new classes arrive, a guarantee that Hoeffding-based criteria cannot provide.
% The analysis applies a per-decision union bound over features and thresholds,
% standard in the streaming tree literature~\cite{domingos2000,derosa2015splitting};
% a time-uniform schedule via Ville-type inequalities is left for future work.
\end{remark}
\subsection{Composed Split Error with Sketched Features}
\label{sec:theory-composed}

\begin{theorem}[Composed Split Error at Fixed Threshold]
\label{thm:composed}
Let $F^*$ be the expected true Gini gain and $\widehat{F}$ the estimate
from a quantile sketch with rank error $\varepsilon_{\mathrm{sketch}}$.
For any $\delta, \delta_{\mathrm{sketch}} \in (0,1)$, with probability
at least $1 - \delta - \delta_{\mathrm{sketch}}$:
\begin{equation}
  |\widehat{F} - F^*| \leq
  \underbrace{\sqrt{\frac{8\ln(2dm/\delta)}{n}}}_{\text{statistical}}
  + \underbrace{4\varepsilon_{\mathrm{sketch}}}_{\text{sketch}}.
  \label{eq:composed-error}
\end{equation}
The statistical term uses factor 8 rather than 32 because bounding
$|\widehat{F} - F^*|$ at a fixed threshold requires only a two-sided
single-quantity concentration argument, not the simultaneous two-tail
control needed by the GapTest in Eq.~\eqref{eq:radius}.
Full proof in Appendix~\ref{app:proof-composed}.
The factor~$4\varepsilon_{\mathrm{sketch}}$ follows from Theorem~\ref{thm:tight}:
rank error shifts routing counts by $n\varepsilon_{\mathrm{sketch}}$, costing
$(4/n)\cdot n\varepsilon_{\mathrm{sketch}}$ in Gini gain.
\end{theorem}

\begin{corollary}[Full Composed Error with Optimization Displacement]
\label{cor:opt-displacement}
Under Gaussianity of $X_j \mid c$ and finite condition number
$\kappa = \sigma_{\max}/\sigma_{\min}$, with probability at least
$1 - \delta - \delta_{\mathrm{sketch}}$:
\begin{equation}
  F(S;\,v^*) - F(S;\,\hat{v}^*)
  \leq
  \sqrt{\frac{8\ln(2dm/\delta)}{n}}
  + 4(1+\kappa)\,\varepsilon_{\mathrm{sketch}},
  \label{eq:composed-full}
\end{equation}
where $\hat{v}^*$ is the sketch-derived threshold. The additional
$4\kappa\,\varepsilon_{\mathrm{sketch}}$ over Theorem~\ref{thm:composed}
accounts for threshold displacement: KLL rank error $\varepsilon_{\mathrm{sketch}}$
implies $|\hat{v}^* - v^*| \leq \varepsilon_{\mathrm{sketch}}\,\sigma_{\max}\sqrt{2\pi}$
under Gaussianity, and by the Lipschitz constant of Lemma~\ref{lem:gini-lipschitz}
($L_v = O(1/\sigma_{\min})$), this propagates to a Gini loss of
$4\kappa\,\varepsilon_{\mathrm{sketch}}$.
Accuracy saturates at $k \geq 32$ and the gap to exact-threshold
evaluation falls below $0.1\%$ (Table~\ref{tab:ablation-k-alpha}),
confirming the displacement term is negligible in practice.
Full proof in Appendix~\ref{app:proof-displacement}.
\end{corollary}
\subsection{Variance Reduction from Bayesian Inheritance}
\label{sec:theory-var}
\begin{proposition}[Post-Split Variance Reduction]
\label{prop:variance}
Let $\alpha \in (0,1)$ be the inheritance discount and
$S_0 = \alpha n_{\mathrm{parent}}$ the inherited pseudo-count total.
\emph{(i) Class probabilities.}
Under the Dirichlet inheritance prior, the posterior variance satisfies
for all $t \geq 0$:
\begin{equation}
  \mathrm{Var}(\theta_c \mid \mathcal{D}_t)
  \leq \tfrac{1}{4(S_0 + 1)}
  = O\!\left(\tfrac{1}{\alpha n_{\mathrm{parent}}}\right).
\end{equation}
\emph{(ii) Feature parameters.}
Under the NIG conjugate prior with $n_0 = \alpha n_c^{\mathrm{parent}}$
pseudo-observations, the marginal posterior of $\mu_j^c$ is
Student-$t$ with $\nu = n_0$ degrees of freedom and scale
$(\sigma_j^c)^2/n_0$. Its variance exists when $n_0 > 2$ and equals:
\begin{equation}
  \mathrm{Var}(\mu_j^c \mid \mathcal{D}_0)
  = \frac{(\sigma_j^c)^2}{n_0} \cdot \frac{\nu}{\nu - 2}
  = \frac{(\sigma_j^c)^2}{\alpha n_c^{\mathrm{parent}} - 2}
  = O\!\left(\tfrac{1}{\alpha n_c^{\mathrm{parent}}}\right).
\end{equation}
For the split feature, $(\sigma_j^c)^2$ is replaced by
$\sigma_{\mathrm{trunc}}^2 \leq (\sigma_j^c)^2$.
Full proof in Appendix~\ref{app:proof-variance}.
\end{proposition}

\begin{remark}
Both posteriors tighten as $O(1/(\alpha n_{\mathrm{parent}}))$, and
since the McDiarmid GapTest maximises $n_{\mathrm{parent}}$ before any
split, conservative splitting and non-zero inheritance act as
complementary mechanisms that jointly minimise post-split uncertainty.
Any $\alpha \geq 0.2$ recovers near-full performance; $\alpha = 0$
replicates the \emph{w/o Inheritance} collapse (Table~\ref{tab:ablation-k-alpha}).
\end{remark}
% -------------------------------------------------------
% EXPERIMENTS
% -------------------------------------------------------
\section{Experiments}
\label{sec:experiments}

\subsection{Experimental Setup}
\label{sec:setup}

\noindent
\textbf{Datasets.}
We evaluate on two benchmark categories. \emph{Standard benchmarks} include
two synthetic Gaussian mixtures (Synth-10, Synth-50) and eight real-world
tabular datasets: Covertype, Split-MNIST, Pendigits, Shuttle, Letter, Wine,
Iris, and HAR. Categorical attributes are one-hot encoded before streaming.
\emph{Stress-test synthetics} (Section~\ref{sec:stress}) isolate specific
failure modes of parametric leaf predictors: Multimodal, Skewed,
AngularSectors, Antipodal, Ring, HeavyTail, NoisyFeature, and ConceptDrift.
For Split-MNIST, raw pixels ($d{=}784$) are reduced to $d{=}50$ via PCA
fitted on the first $1{,}000$ task-0 samples only. The frozen projection is
then applied uniformly to all methods.

\noindent
\textbf{Baselines.}
We compare two groups of methods. \emph{Tree-based:} \vfdt{}, Rutkowski,
EFDT, HAT, and \arf{}-NR (10 trees, drift detection disabled). EFDT and HAT
under default drift-reset collapse under strict Online CIL
(Appendix~\ref{app:noreset}); we report their no-reset variants in the main
comparison only. \emph{Non-tree:} \slda{}, \sqda{}, NCM, RLS. We also compare against PEC~\cite{zajkac2024}, a per-class exemplar-free baseline that avoids global coupling via independent class models. All
methods follow the strict Online CIL protocol: single-pass, no raw data
retention, sequentially expanding $K$\footnote{Note that these are the SOTA methods for our setting: online CIL with raw tabular features. More recent, technically heavily involved, CIL methods require pretrained backbones so they do not apply directly to our setting~\cite{goswami2023fecam, zhuang2024gacl, petit2023fetril}.}. 

\noindent
\textbf{Hyperparameters.}
Sketch capacity $k$, inheritance discount $\alpha$, failure probability
$\delta$, grace period, and threshold generation policy are fixed globally
without per-dataset tuning; defaults are $k{=}64$, $\alpha{=}0.6$,
$\delta{=}0.10$, grace${=}200$ (full configurations in
Table~\ref{tab:hyperparams}). Sensitivity to $k$, $\alpha$, and
$\delta$ is in Table~\ref{tab:ablation-k-alpha}; sensitivity of
\mist{}-K to $\epsilon_s$ and $\beta$ is in
Appendix~\ref{app:mistk-sensitivity}.

\begin{table}[h]
\caption{Final mean accuracy among \textbf{tree-based methods} (mean $\pm$
std over 5 seeds). \textbf{Bold}: best within group. All methods follow the
no-reset CIL-adapted protocol; default drift-reset results for EFDT and HAT
are reported in Appendix~\ref{app:noreset}.
\arf{}-NR: Adaptive Random Forest with drift detection disabled (10 trees).}
\label{tab:tree-results}
\centering
\resizebox{\textwidth}{!}{%
\setlength{\tabcolsep}{3pt}
\begin{tabular}{lccccccccccc}
\toprule
\textbf{Method}
  & Synth-10 & Synth-50 & Wine & Iris
  & Covertype & Split-MNIST & Pendigits & Shuttle & Letter & HAR & \textbf{Mean} \\
\midrule
\mist{}-G (ours)
  & 0.994$\pm$0.004 & 0.999$\pm$0.002 & \textbf{0.970}$\pm$0.020
  & \textbf{0.933}$\pm$0.030 & 0.559$\pm$0.013 & \textbf{0.863}$\pm$0.003
  & \textbf{0.889}$\pm$0.023 & 0.733$\pm$0.033
  & 0.678$\pm$0.005 & 0.607$\pm$0.013 & 0.8225 \\
\mist{}-K (ours)
  & \textbf{0.996}$\pm$0.002 & \textbf{1.000}$\pm$0.001 & 0.957$\pm$0.038
  & 0.927$\pm$0.013 & \textbf{0.694}$\pm$0.006 & 0.772$\pm$0.007
  & 0.877$\pm$0.020 & \textbf{0.783}$\pm$0.067
  & \textbf{0.692}$\pm$0.009 & \textbf{0.625}$\pm$0.009 & \textbf{0.8323} \\
\midrule
\vfdt{}
  & 0.998$\pm$0.002 & 0.396$\pm$0.003 & \textbf{0.970}$\pm$0.020
  & \textbf{0.933}$\pm$0.030 & \textbf{0.605}$\pm$0.010 & \textbf{0.874}$\pm$0.004
  & 0.508$\pm$0.039 & 0.417$\pm$0.145
  & 0.372$\pm$0.056 & 0.324$\pm$0.008 & 0.6397 \\
Rutkowski
  & \textbf{1.000}$\pm$0.000 & \textbf{1.000}$\pm$0.000 & \textbf{0.970}$\pm$0.020
  & \textbf{0.933}$\pm$0.030 & 0.559$\pm$0.013 & 0.851$\pm$0.003
  & \textbf{0.863}$\pm$0.010 & \textbf{0.527}$\pm$0.099
  & \textbf{0.678}$\pm$0.005 & \textbf{0.614}$\pm$0.018 & \textbf{0.7995} \\
EFDT (NR)
  & 0.219$\pm$0.038 & 0.040$\pm$0.000
  & \textbf{0.970}$\pm$0.020 & \textbf{0.933}$\pm$0.030
  & 0.277$\pm$0.002 & 0.182$\pm$0.009 & 0.272$\pm$0.057 & 0.143$\pm$0.000
  & 0.168$\pm$0.026 & 0.202$\pm$0.066 & 0.3406 \\
HAT (NR)
  & 0.423$\pm$0.049 & 0.084$\pm$0.017
  & \textbf{0.970}$\pm$0.020 & \textbf{0.933}$\pm$0.030
  & 0.262$\pm$0.032 & 0.324$\pm$0.037 & 0.614$\pm$0.037 & 0.220$\pm$0.063
  & 0.237$\pm$0.032 & 0.307$\pm$0.054 & 0.4374 \\
\arf{}-NR
  & \textbf{1.000}$\pm$0.000 & \textbf{1.000}$\pm$0.000
  & 0.952$\pm$0.010 & 0.927$\pm$0.025
  & 0.458$\pm$0.013 & 0.782$\pm$0.009 & 0.846$\pm$0.009 & 0.444$\pm$0.029
  & 0.568$\pm$0.012 & 0.601$\pm$0.010 & 0.7578 \\
\bottomrule
\end{tabular}}
\end{table}

\begin{table}[h]
\caption{Final mean accuracy among \textbf{non-tree baselines} (mean $\pm$
std over 5 seeds), included as reference. \textbf{Bold}: best within group.
These methods maintain a single global representation across all classes and
do not share the structural modularity properties of streaming decision trees;
the comparison is included to characterise the accuracy regime instead of showing a direct competition.}
\label{tab:nontree-results}
\centering
\resizebox{\textwidth}{!}{%
\setlength{\tabcolsep}{3pt}
\begin{tabular}{lccccccccccc}
\toprule
\textbf{Method}
  & Synth-10 & Synth-50 & Wine & Iris
  & Covertype & Split-MNIST & Pendigits & Shuttle & Letter & HAR & \textbf{Mean}\\
\midrule
\mist{}-G (ours)  & 0.994$\pm$0.004 & 0.999$\pm$0.002 & \textbf{0.970}$\pm$0.020
  & \textbf{0.933}$\pm$0.030 & 0.559$\pm$0.013 & \textbf{0.863}$\pm$0.003
  & \textbf{0.889}$\pm$0.023 & 0.733$\pm$0.033
  & 0.678$\pm$0.005 & 0.607$\pm$0.013 & 0.8225 \\
\mist{}-K (ours)
  & \textbf{0.996}$\pm$0.002 & \textbf{1.000}$\pm$0.001 & 0.957$\pm$0.038
  & 0.927$\pm$0.013 & \textbf{0.694}$\pm$0.006 & 0.772$\pm$0.007
  & 0.877$\pm$0.020 & \textbf{0.783}$\pm$0.067
  & \textbf{0.692}$\pm$0.009 & \textbf{0.625}$\pm$0.009 & \textbf{0.8323} \\
\midrule
\slda{}
  & \textbf{1.000}$\pm$0.000 & \textbf{1.000}$\pm$0.000
  & 0.979$\pm$0.021 & \textbf{0.973}$\pm$0.039
  & \textbf{0.673}$\pm$0.007 & 0.854$\pm$0.005 & 0.879$\pm$0.007 & 0.529$\pm$0.000
  & 0.718$\pm$0.004 & \textbf{0.815}$\pm$0.001 & 0.8420 \\
\sqda{}
  & \textbf{1.000}$\pm$0.000 & 0.999$\pm$0.002
  & \textbf{0.987}$\pm$0.016 & 0.967$\pm$0.030
  & 0.528$\pm$0.004 & \textbf{0.940}$\pm$0.003
  & \textbf{0.982}$\pm$0.002 & \textbf{0.774}$\pm$0.029
  & \textbf{0.895}$\pm$0.006 & 0.791$\pm$0.005 & \textbf{0.8863} \\
NCM
  & \textbf{1.000}$\pm$0.000 & \textbf{1.000}$\pm$0.000
  & 0.981$\pm$0.010 & 0.860$\pm$0.039
  & 0.603$\pm$0.012 & 0.794$\pm$0.003 & 0.817$\pm$0.009
  & 0.549$\pm$0.030 & 0.602$\pm$0.004 & 0.639$\pm$0.007 & 0.7845 \\
RLS
  & 0.999$\pm$0.002 & \textbf{1.000}$\pm$0.000
  & 0.981$\pm$0.028 & 0.827$\pm$0.049
  & 0.659$\pm$0.007 & 0.818$\pm$0.006 & 0.857$\pm$0.008 & 0.350$\pm$0.057
  & 0.596$\pm$0.004 & \textbf{0.815}$\pm$0.001 & 0.7902 \\
PEC
  & 1.000$\pm$0.001 & 0.952$\pm$0.018 & 0.782$\pm$0.057 & 0.820$\pm$0.040 & 0.650$\pm$0.007 & 0.898$\pm$0.009 & 0.961$\pm$0.002 & 0.421$\pm$0.005 & 0.760$\pm$0.005 & 0.614$\pm$0.011 & 0.7857 \\
\bottomrule
\end{tabular}}
\end{table}

\subsection{Main Results}
\label{sec:results}
Among tree-based methods (Table~\ref{tab:tree-results}), \mist{}-G and
\mist{}-K are the strongest single-tree configurations across the benchmark
suite. \efdt{}~(NR) collapses on most benchmarks, consistent with
cold-start instability compounding under more aggressive split triggers
without inheritance. \arf{}-NR, with 10 base trees, falls below both
\mist{} configurations on most real-stream benchmarks 
with memory overhead that depends on the dataset regime
(Table~\ref{tab:memory}, Appendix~\ref{app:memory-full}).
% while requiring
% substantially more memory (Table~\ref{tab:memory}, Appendix~\ref{app:memory-full}).
% with memory efficiency that varies by regime: ARF-NR uses more memory
% on low-$K$ datasets but less on higher-$K$ ones where per-class KLL
% sketches dominate (Table~\ref{tab:memory}, Appendix~\ref{app:memory-full}).

On near-Gaussian, low-$K$ streams (Synth-10, Wine, Iris), \vfdt{} matches
\mist{}-G on final accuracy, reflecting an honest trade-off: its more
aggressive Hoeffding criterion yields faster structural adaptation when class
geometry is simple. The forgetting results (Appendix~\ref{app:forgetting-full})
provide the complementary picture, showing \mist{}-G achieving consistently lower
forgetting, confirming that accuracy parity on these streams comes at the
cost of greater instability under class arrival. 
On Letter ($K{=}26$, $d{=}16$), zero splits under default settings
reflects the GapTest withholding commitment; with $K{=}26$ classes
and $d{=}16$ features the Gini gain gap between candidate splits
does not exceed the McDiarmid radius, indicating the root partition
is insufficiently informative to warrant a split under the exemplar-free
constraint.
%
Table~\ref{tab:nontree-results} reports non-tree baselines as an accuracy
reference. Observe that our single streaming tree with the proposed calibrations is
competitive with global parametric methods on near-Gaussian benchmarks. This
confirms that local updates and bounded forgetting need not sacrifice
accuracy in this regime. The picture reverses on non-Gaussian geometry
(Section~\ref{sec:stress}), where global methods collapse. Average forgetting
results and prequential accuracy curves are in
Appendices~\ref{app:prequential}--\ref{app:forgetting-full}.
\subsection{Ablation Study}
\begin{table}[h]
\caption{Core and split-criterion ablations (average accuracy over 
ten standard datasets, mean over 5 seeds).}
\label{tab:ablation-core}
\centering
\small
\setlength{\tabcolsep}{3pt}
\begin{tabular}{lclc}
\toprule
Variant & Avg Acc. & Description & $\Delta$ \\
\midrule
\multicolumn{4}{l}{\textit{Full architecture}} \\
\mist{}-G (default) & 0.8225 & McDiarmid + sketch + inheritance + GNB leaf & --- \\
\midrule
\multicolumn{4}{l}{\textit{Architectural variants}} \\
w/o Inheritance    & 0.7439 & No parent-to-child transfer         & $-$\textbf{9.55}\%  \\
w/o Sketching      & 0.8204 & Gaussian threshold approximation    & $-$0.26\%  \\
w/ Hoeffding       & 0.8309 & Hoeffding split rule                & $+$1.02\%  \\
Majority-vote leaf & 0.2306 & No density model                    & $-$\textbf{71.97}\% \\
w/ SampleSplit     & 0.8227 & Sample-based threshold generation   & $+$0.02\%  \\
w/ DenseThreshold  & 0.8218 & Dense threshold grid                & $-$0.09\%  \\
\midrule
\multicolumn{4}{l}{\textit{Split criterion variants}} \\
GapTest (default)    & 0.8225 & Gap between best and second-best   & ---        \\
GainTest             & 0.8324 & Best gain vs.\ radius              & $+$1.20\%  \\
GainTestNoUnion      & 0.8405 & GainTest without union-bound term  & $+$2.19\%  \\
GaussianInheritance  & 0.8228 & GapTest + Gaussian CDF projection  & $+$0.04\%  \\
\bottomrule
\end{tabular}
\end{table}
% \subsection{Ablation Study}
% \begin{table}[h]
% \caption{Core and split-criterion ablations (average accuracy over 
% ten standard datasets, mean over 5 seeds).}
% \label{tab:ablation-core}
% \centering
% \small
% \setlength{\tabcolsep}{3pt}
% \begin{tabular}{lclc}
% \toprule
% Variant & Avg Acc. & Description & $\Delta$ \\
% \midrule
% \multicolumn{4}{l}{\textit{Full architecture}} \\
% \mist{}-G (default) & 0.8655 & McDiarmid + sketch + inheritance + GNB leaf & --- \\
% \midrule
% \multicolumn{4}{l}{\textit{Architectural variants}} \\
% w/o Inheritance    & 0.7663 & No parent-to-child transfer         & $-$11.46\% \\
% w/o Sketching      & 0.8663 & Gaussian threshold approximation    & $+$0.09\%  \\
% w/ Hoeffding       & 0.8656 & Hoeffding split rule                & $+$0.01\%  \\
% Majority-vote leaf & 0.3241 & No density model                    & $-$62.55\% \\
% w/ SampleSplit     & 0.8651 & Sample-based threshold generation   & $-$0.05\%  \\
% w/ DenseThreshold  & 0.8637 & Dense threshold grid                & $-$0.21\%  \\
% \midrule
% \multicolumn{4}{l}{\textit{Split criterion variants}} \\
% GapTest (default)    & 0.8655 & Gap between best and second-best   & --- \\
% GainTest             & 0.8611 & Best gain vs.\ radius              & $-$0.51\%  \\
% GainTestNoUnion      & 0.8639 & GainTest without union-bound term  & $-$0.18\%  \\
% GaussianInheritance  & 0.8660 & GapTest + Gaussian CDF projection  & $+$0.06\%  \\
% \bottomrule
% \end{tabular}
% \end{table}

We focus on two main ablation tests. Removing inheritance causes a relative drop of
$9.55\%$ ($7.86$ pp absolute), confirming that post-split cold-start
instability is the primary failure mode in the exemplar-free setting.
Replacing GNB leaves with majority-vote causes a $71.97\%$ relative
collapse ($59.19$ pp absolute): without class-conditional likelihood
estimates, the tree cannot distinguish between classes sharing the same
partition cell.
The McDiarmid versus Hoeffding comparison shows negligible accuracy difference
($+1.02\%$) on near-Gaussian benchmarks. The structurally meaningful difference is 
in Figure~\ref{fig:tree-size}, 
Appendix~\ref{app:tree-dynamics}): Hoeffding triggers substantially
more splits as $K_{\mathrm{local}}$ grows, while McDiarmid maintains compact
trees. 
GainTest yields marginal accuracy gains ($+1.20\%$) at the cost of
substantially larger trees (Figure~\ref{fig:tree-size}), confirming the
McDiarmid radius as a structural regulariser.
Accuracy is robust to $k$, $\alpha$, and $\delta$ over broad ranges
(Appendix~\ref{app:sensitivity}): performance saturates at $k \geq 32$ and
$\alpha \geq 0.2$, confirming that \mist{} can be deployed at low memory
budgets with minimal accuracy cost. Setting $\alpha = 0$ collapses to
$0.7439$, exactly matching the \emph{w/o Inheritance} ablation.

\subsection{Stress Test on Non-Gaussian Geometry}
\label{sec:stress}

\begin{table}[h]
\caption{Stress-test results (mean $\pm$ std over 5 seeds) on non-Gaussian
synthetic streams.}
\label{tab:stress}
\centering
\resizebox{\textwidth}{!}{%
\setlength{\tabcolsep}{3pt}
\begin{tabular}{lccccccccc}
\toprule
Method & Antipodal & Multimodal & Ring & Skewed & HeavyTail
       & NoisyFeature & ConceptDrift & AngularSectors & \textbf{Mean} \\
\midrule
\mist{}-G & 0.834$\pm$0.033 & 0.919$\pm$0.032 & 0.975$\pm$0.019
          & 0.899$\pm$0.012 & 0.881$\pm$0.014 & 0.808$\pm$0.106
          & \textbf{0.993}$\pm$0.012 & 0.241$\pm$0.010 & 0.8188 \\
\mist{}-K & 0.829$\pm$0.021 & 0.983$\pm$0.005 & 0.968$\pm$0.026
          & 0.978$\pm$0.009 & 0.897$\pm$0.022 & 0.768$\pm$0.121
          & 0.988$\pm$0.011 & \textbf{0.436}$\pm$0.017 & \textbf{0.8559} \\
\midrule
\slda{}   & 0.171$\pm$0.042 & 0.756$\pm$0.053 & 0.978$\pm$0.020
          & 0.897$\pm$0.012 & \textbf{0.898}$\pm$0.025 & 0.811$\pm$0.104
          & \textbf{0.993}$\pm$0.012 & 0.240$\pm$0.041 & 0.7180 \\
\sqda{}   & \textbf{0.984}$\pm$0.007 & \textbf{0.999}$\pm$0.001
          & 0.975$\pm$0.021 & \textbf{1.000}$\pm$0.000
          & 0.858$\pm$0.012 & 0.798$\pm$0.109
          & \textbf{0.993}$\pm$0.012 & 0.235$\pm$0.015 & 0.8553 \\
\arf{}-NR & 0.862$\pm$0.025 & 0.947$\pm$0.018 & \textbf{0.979}$\pm$0.012
          & 0.921$\pm$0.016 & 0.893$\pm$0.019 & \textbf{0.832}$\pm$0.091
          & 0.991$\pm$0.009 & 0.255$\pm$0.021 & 0.8350 \\
PEC       & 0.918$\pm$0.010 & 0.964$\pm$0.017 & 0.926$\pm$0.035 & 0.972$\pm$0.009 & 0.801$\pm$0.012 & 0.233$\pm$0.025 & 0.984$\pm$0.018 & 0.302$\pm$0.046 & 0.7625 \\
\bottomrule
\end{tabular}}
\end{table}
Recall that in Table~\ref{tab:nontree-results} non-tree benchmarks
outperform our solutions due to their utilisation of the Gaussian
behaviour of the data. To test how robust these methods are under
non-Gaussian geometry, we run additional stress tests as follows:
On AngularSectors, \mist{}-K is the only configuration that substantially
exceeds chance: \sqda{} ($0.235$) and \arf{}-NR ($0.255$) both collapse
because ensemble diversity cannot compensate for a shared Gaussian leaf
assumption---the failure is a bias problem, not a variance problem. On
Multimodal and Skewed, \mist{}-K substantially outperforms \mist{}-G and
approaches \sqda{}, confirming the sketch captures non-Gaussian structure
without raw-data storage. The failure of \slda{} on Antipodal
while \sqda{} succeeds illustrates the importance of per-class
covariance in non-spherical settings. Figure~\ref{fig:stress-vis} provides
PCA visualisations of all stress-test streams.

\mist{} carries higher memory than lightweight baselines due to per-leaf,
per-class quantile sketches. Memory scales as $O(dkL_tK_t)$ worst-case,
but most leaves quickly satisfy $K_{\mathrm{local}}^{(\ell)} \ll K_t$ as
leaves purify (Appendix~\ref{app:memory-scaling}). Among Pareto-optimal
configurations on Covertype, \mist{}-K achieves the highest accuracy;
\slda{} offers near-best accuracy at low memory; \mist{}-G uniquely
provides formal split-calibration guarantees and non-Gaussian robustness
where global methods collapse. Sketch-quantile initialisation matches
truncated-Gaussian accuracy on near-Gaussian data but is $1.27\times$
slower, confirming inheritance's benefit lies in the warm-start mechanism;
on outlier-contaminated data it yields gains up to $+33.5$~pp
(Appendix~\ref{app:sketch-vs-gaussian}).

% -------------------------------------------------------
% CONCLUSION
% Tone: concise, honest about limitations, forward-looking
%   toward Prof's modular vision
% No overclaiming - just what was shown
% -------------------------------------------------------

\vspace{-0.2cm}
\section{Conclusion}
\label{sec:conclusion}
Streaming decision tree failure in online CIL traces to two coupled sources: a split criterion that degrades as $K$ grows, and the absence of parent-to-child knowledge transfer. \mist{} resolves this through three integrated components: a tight, $K$-independent McDiarmid confidence radius ($c_i{=}4/n$) for Gini splitting that prevents premature splits; a Bayesian inheritance mechanism for post-split stabilization; and unified KLL quantile sketches for continuous split evaluation and geometry-adaptive prediction.
%
Empirically, \mist{}-G is highly competitive with global parametric methods on near-Gaussian streams. On strongly non-Gaussian geometry, \mist{}-K is uniquely robust, succeeding where per-class discriminants and 10-tree ensembles collapse to near-chance (e.g., AngularSectors). This regime-dependence explicitly validates our dual parametric/non-parametric design.

Regarding limitations, \mist{}-G's projection inherently relies on Gaussian assumptions, which \mist{}-K bypasses via sketching; extending this infrastructure with copulas could further capture feature dependencies. Crucially, deriving time-uniform confidence sequences via nonnegative supermartingales \citep{Howard2021} is a major focus for future work. This would replace our current per-decision union bound, formally eliminating multiple-testing effects over infinite horizons and further optimizing split timing.

Ultimately, \mist{} demonstrates that streaming decision trees provide a naturally modular architecture for open-world continual learning. By strictly decoupling local predictive modules while enabling statistically grounded knowledge transfer at split time, structured partitioning offers a principled foundation for learners that grow without catastrophic forgetting.

\newpage
\bibliographystyle{plain}
\bibliography{references}

\newpage
\appendix
% ============================================================
%  Appendix: Related Work on Online Continual Learning
%  For inclusion in MIST (NeurIPS 2026) appendix
% ============================================================



\section{Broader Context: Online Continual Learning}
\label{app:related_work}

\subsection{Catastrophic forgetting and the continual-learning taxonomy}
\label{app:related_work:taxonomy}

The tendency of connectionist learners to overwrite previously acquired knowledge was
formalised as \emph{catastrophic interference} by~\cite{mccloskey1989catastrophic}
and~\cite{ratcliff1990connectionist}, and revisited for modern networks
by~\cite{french1999catastrophic} and~\cite{goodfellow2013empirical}.
Modern continual learning (CL) is commonly organised along three scenarios, including \textbf{task-incremental, domain-incremental, and class-incremental}, formalised
by~\cite{vandeven2019three} and later refined in~\cite{vandeven2022three};
the categorisation was independently validated by~\cite{hsu2018reevaluating}.
Class-incremental learning (CIL) is widely regarded as the hardest setting because
the task identifier is not supplied at inference and the classifier must discriminate
between all classes ever seen, requiring calibrated cross-task decision boundaries.
Broader surveys include~\cite{parisi2019continual,delange2022survey}, and,
for the specifically \emph{online} (single-pass, small-batch) regime,~\cite{mai2022online}.
Streaming settings aggravate forgetting relative to the batch case because each sample
is seen only once and stationarity cannot be assumed, motivating drift-aware
evaluation~\cite{ghunaim2023realtime}.

\subsection{Regularisation-based methods}
\label{app:related_work:regularisation}

Regularisation approaches anchor parameters deemed important for past tasks.
Elastic Weight Consolidation (EWC) uses a Fisher-information
penalty~\cite{kirkpatrick2017overcoming}; Synaptic Intelligence accumulates path
integrals of parameter sensitivity~\cite{zenke2017continual}; Memory-Aware Synapses
uses gradient magnitudes of the learned function~\cite{aljundi2018memory};
Learning without Forgetting employs knowledge distillation from the previous model
as a functional prior~\cite{li2018learning}.
Online variants such as Progress \& Compress~\cite{schwarz2018progress} and
RWalk~\cite{chaudhry2018riemannian} were developed to amortise Fisher estimation
over streams.
These methods are effective for short task sequences but degrade sharply in unbounded
class-incremental streams: importance estimates drift, the quadratic penalty cannot
simultaneously protect arbitrarily many past classes, and they do not address the
representation--classifier mismatch at the core of CIL.

\subsection{Replay and rehearsal}
\label{app:related_work:replay}

Replay stores or synthesises past data to counteract forgetting.
Experience Replay~\cite{rolnick2019experience}, iCaRL~\cite{rebuffi2017icarl},
GEM~\cite{lopezpaz2017gem} and its efficient variant A-GEM~\cite{chaudhry2019agem}
constrain updates with respect to a memory buffer; tiny-memory baselines are analysed
in~\cite{chaudhry2019tiny}.
DER++~\cite{buzzega2020der} augments replay with logit distillation,
while MIR~\cite{aljundi2019mir} and GSS~\cite{aljundi2019gss} target the
\emph{online} regime with interference-aware retrieval and gradient-based sample selection.
ER-ACE~\cite{caccia2022eracer} mitigates representation drift caused by
na\"{i}ve cross-entropy replay.
Generative replay replaces the buffer with a learned generator~\cite{shin2017dgr},
and GDumb~\cite{prabhu2020gdumb} shows that a buffer-trained model is a strong baseline.
Selective replay for lifelong reinforcement learning was proposed
by~\cite{isele2018selective}.
All of these methods \textbf{violate the exemplar-free constraint} relevant to
privacy-sensitive tabular applications such as finance and healthcare, and generative
replay additionally compounds forgetting via generator drift.

\subsection{Architecture-expansion and parameter isolation}
\label{app:related_work:architecture}

A complementary family grows or partitions the network per task:
Progressive Networks~\cite{rusu2016progressive},
PackNet~\cite{mallya2018packnet},
Dynamically Expandable Networks~\cite{yoon2018lifelong},
HAT~\cite{BifetGavalda2009HAT},
Piggyback~\cite{mallya2018piggyback}, and
Expert Gate~\cite{aljundi2017expert}.
These methods largely avoid forgetting by construction but require a task
identifier at inference or introduce unbounded parameter growth, which is
incompatible with an open-ended class stream.

\subsection{Exemplar-free deep continual learning}
\label{app:related_work:exemplar_free}

Exemplar-free CIL methods retain only compact statistics.
PASS~\cite{zhu2021prototype} uses prototype augmentation and self-supervision;
IL2M~\cite{belouadah2019il2m} rebalances logits with class statistics;
SSRE~\cite{zhu2022self} expands representations without exemplars;
PRAKA~\cite{shi2023prototype} introduces prototype reminiscence and asymmetric
knowledge aggregation;
FeTrIL~\cite{petit2023fetril} translates old-class features from new-class geometry;
FeCAM~\cite{goswami2023fecam} exploits heterogeneous class covariances.
Deep SLDA~\cite{hayes2020lifelong} freezes a backbone and maintains a
streaming Gaussian classifier.
Prompt-based methods such as L2P~\cite{wang2022l2p},
DualPrompt~\cite{wang2022dualprompt}, and CODA-Prompt~\cite{smith2023coda}
exploit large pre-trained vision transformers with learnable prompt pools; they
depend on the quality and transferability of the frozen backbone, a resource
unavailable in most tabular domains.

\subsection{Online and streaming-specific methods}
\label{app:related_work:online}

REMIND~\cite{hayes2020remind} and Latent Replay~\cite{pellegrini2020latent}
compress activations to enable single-pass rehearsal;
\cite{lesort2020continual} surveys the online robotics setting;
\cite{ghunaim2023realtime} argues for throughput-aware evaluation.
In contrast to multi-epoch offline CIL, these works highlight the stringent
per-sample compute budget that governs truly online learners.

\subsection{Tabular continual learning and streaming decision trees}
\label{app:related_work:tabular}

Continual learning on \textbf{tabular} streams is markedly under-explored.
Recent work considers tabular CL in energy-constrained~\cite{wang2025imlp}
and contrastive or out-of-distribution settings~\cite{tccl2025}, as well as
financial-audit anomaly detection~\cite{hemati2022continual}; a pseudo-rehearsal
approach is reported in~\cite{tril3_2025}.
These efforts remain either replay-based or confined to narrow domains, leaving
class-incremental tabular learning largely open.

The complementary data-stream-mining literature offers decision-tree learners with
provable split guarantees: the Hoeffding-bound VFDT~\cite{DomingosHulten2000} and
its concept-drift extension CVFDT~\cite{HultenSpencerDomingos2001}; the Extremely
Fast Decision Tree (EFDT)~\cite{ManapragadaWebbSalehi2018}; the Hoeffding Adaptive
Tree~\cite{BifetGavalda2009HAT}; streaming
ensembles~\cite{BifetGavalda2009HAT,gomes2017};
and change-detection mechanisms ADWIN~\cite{BifetGavalda2007ADWIN} and
DDM~\cite{GamaMedasCastilloRodrigues2004}.
Crucially, \cite{rutkowski2013} show that VFDT's use of
Hoeffding's inequality is statistically inappropriate for information-gain-type split
criteria and replace it with the tighter \textbf{McDiarmid bound}.
However, none of these streaming-tree methods are designed or analysed for
class-incremental learning: new classes invalidate cached sufficient statistics at
existing splits, pure leaf counts drift as the label space grows, and classical drift
detectors conflate concept drift with class emergence.

\subsection{Positioning of MIST}
\label{app:related_work:positioning}

MIST sits at the intersection of the gaps identified above.
It is \textbf{exemplar-free} (no raw sample buffers or generative models,
unlike~\S\ref{app:related_work:replay}),
\textbf{single-pass and online} (unlike most methods
in~\S\ref{app:related_work:regularisation},
\S\ref{app:related_work:architecture}, and \S\ref{app:related_work:exemplar_free}),
\textbf{class-set-unbounded} (unlike parameter-isolation methods
in~\S\ref{app:related_work:architecture} and short-horizon regularisers
in~\S\ref{app:related_work:regularisation}), and \textbf{designed for tabular streams}
(unlike methods in~\S\ref{app:related_work:online}--\S\ref{app:related_work:tabular},
which presuppose visual pre-training or fixed feature spaces).

Methodologically, MIST replaces the statistically loose Hoeffding split test of
VFDT and EFDT with a \textbf{McDiarmid-calibrated criterion} in the spirit
of~\cite{rutkowski2013}, and extends it with \textbf{Bayesian inheritance of
class-conditional leaf statistics} so that newly emerging classes borrow strength
from ancestor nodes rather than restarting estimation.
This yields a streaming decision-tree learner that tolerates open class sets
without replay, capacity expansion, or a pre-trained backbone.
A detailed analysis of the GapTest threshold and the tightness of the
McDiarmid constant $c_i = 4/n$ is provided in Section~\ref{sec:theory}
and Appendix~\ref{app:proofs}.
% ============================================================
%  MIST — Full Supplementary Material
%  Structure:
%    A  Algorithm Pseudocode
%    B  Extended Experimental Setup
%    C  Main Results: Extended Analysis
%    D  Ablation and Sensitivity Analysis
%    E  Memory and Computational Analysis
%    F  Theoretical Empirical Validation
%    G  Stress-Test Dataset Descriptions
%    H  Mathematical Proofs
% ============================================================

% \appendix

% ============================================================
\section{Algorithm Pseudocode}
\label{app:algorithms}
% ============================================================

Algorithm~\ref{alg:mist} in the main text gives the high-level MIST update and
predict loop. We provide here the two subroutines that it invokes:
Algorithm~\ref{alg:split} details the Bayesian inheritance protocol executed at
every split event, and Algorithm~\ref{alg:mistk} details the sketch-based leaf
likelihood used by \mist{}-K. Both pseudocode listings use the notation
established in Section~4 of the main paper; $\phi$ and $\Phi$ denote the
standard normal PDF and CDF respectively.

\begin{algorithm}[h]
\caption{\mist{}: Split with Bayesian Inheritance}
\label{alg:split}
\begin{algorithmic}[1]
\REQUIRE Leaf $\ell$, split feature $j^*$, threshold $v^*$,
         discount $\alpha \in (0,1]$
\FOR{each child $s \in \{\mathrm{left}, \mathrm{right}\}$}
  \FOR{each class $c$}
    \IF{$j^*$ is continuous}
      \STATE $\zeta \leftarrow (v^* - \mu_{j^*}^c)/\sigma_{j^*}^c$
      \STATE $\tilde\Phi \leftarrow \Phi(\zeta)$ if $s{=}\mathrm{left}$,
             else $1{-}\Phi(\zeta)$
      \STATE $\mu_{j^*}^{s,c} \leftarrow \mu_{j^*}^c -
             \sigma_{j^*}^c\,\phi(\zeta)/\tilde\Phi$ if $s{=}\mathrm{left}$,
             else $\mu_{j^*}^c + \sigma_{j^*}^c\,\phi(\zeta)/\tilde\Phi$
             \COMMENT{Eq.~\eqref{eq:mu_left_right}}
      \IF{$s = \mathrm{left}$}
        \STATE $(\sigma_{j^*}^{s,c})^2 \leftarrow \sigma_{j^*}^{c\,2}
               \!\left[1 - \dfrac{\zeta\phi(\zeta)}{\tilde\Phi}
               - \left(\dfrac{\phi(\zeta)}{\tilde\Phi}\right)^{\!2}\right]$
               \COMMENT{upper-truncated; Appendix~\ref{app:truncgauss}}
      \ELSE
        \STATE $(\sigma_{j^*}^{s,c})^2 \leftarrow \sigma_{j^*}^{c\,2}
               \!\left[1 + \dfrac{\zeta\phi(\zeta)}{\tilde\Phi}
               - \left(\dfrac{\phi(\zeta)}{\tilde\Phi}\right)^{\!2}\right]$
               \COMMENT{lower-truncated; sign of linear term is $+$;
               Appendix~\ref{app:truncgauss}}
      \ENDIF
      \STATE $n_c^s \leftarrow \tilde\Phi \cdot \alpha \cdot n_c^\ell$
    \ELSE
      \STATE \COMMENT{categorical one-hot: split feature is deterministic in each child}
      \STATE $v_{\mathrm{cat}} \leftarrow 0$ if $s{=}\mathrm{left}$, else $1$
      \STATE $\mu_{j^*}^{s,c} \leftarrow v_{\mathrm{cat}}$;\quad
             $(\sigma_{j^*}^{s,c})^2 \leftarrow 0$
             \COMMENT{$x_{j^*}{=}v_{\mathrm{cat}}$ w.p.\ 1 in child $s$; inherit no uncertainty}
      \STATE $n_c^s \leftarrow \alpha \cdot n_c(v_{\mathrm{cat}})$,\;
             $n_c(v_{\mathrm{cat}}){=}\#\{i:y_i{=}c,\,x_i^{(\mathrm{attr})}{=}v_{\mathrm{cat}}\}$
    \ENDIF
    \FOR{$k \neq j^*$}
      \STATE $(\mu_k^{s,c},\,\sigma_k^{s,c}) \leftarrow (\mu_k^c,\,\sigma_k^c)$
    \ENDFOR
  \ENDFOR
  \STATE Initialise KLL sketches at $s$ as empty
  \COMMENT{\mist{}-G: unaffected (Gaussian accumulators inherited above).
           \mist{}-K: sketch-based feature likelihoods unavailable until
           new samples arrive; class-probability priors remain valid.}
\ENDFOR
\STATE Remove $\ell$ from active leaves
\end{algorithmic}
\end{algorithm}

\begin{algorithm}[h]
\caption{\mist{}-K: Sketch-Based Leaf Likelihood}
\label{alg:mistk}
\begin{algorithmic}[1]
\REQUIRE Observation $x$, leaf $\ell$, smoothing $\epsilon_s$ (default $1.0$),
         bandwidth multiplier $\beta$,
         Dirichlet pseudo-counts $\{\alpha_c^{(\ell)}\}$ where
         $\alpha_c^{(\ell)} = n_c^{(\ell)} + \epsilon_s$,\\
         $\mathrm{IQS}_{\ell,c}(x_j) = \widehat{F}_{j,c}^{-1}(\min(r+h_0,1))
         - \widehat{F}_{j,c}^{-1}(\max(r-h_0,0))$
         where $r = \widehat{F}_{j,c}(x_j)$ is the sketch-estimated rank
         of $x_j$ and $h_0=0.25$ is the fixed half-window;
         returns $\epsilon_s$ if fewer than two distinct quantiles exist
\FOR{each class $c$}
  \STATE $\log\hat{p}(c) \leftarrow \log\alpha_c^{(\ell)} -
         \log\sum_{c'}\alpha_{c'}^{(\ell)}$
         \COMMENT{Dirichlet prior}
  \FOR{each feature $j$}
    \IF{$j$ is categorical}
      \STATE $\hat{p}_j \leftarrow (n_c^{(\ell)}(j{=}x_j) + \epsilon_s)\;/\;
             (n_c^{(\ell)} + 2\epsilon_s)$
             \COMMENT{Bernoulli likelihood from Dirichlet pseudo-counts;
             $n_c^{(\ell)}(j{=}x_j)$: count of class-$c$ samples with
             $x_j$ matching the observed category}
      \STATE $\log\hat{p}(x_j \mid c) \leftarrow \log\hat{p}_j$
    \ELSE
      \STATE $h_j \leftarrow \beta \cdot \mathrm{IQS}_{\ell,c}(x_j)$
             \COMMENT{local inter-quantile spacing}
      \STATE $\hat{F}^+ \leftarrow \widehat{F}_{j,c}(x_j + h_j)$;\quad
             $\hat{F}^- \leftarrow \widehat{F}_{j,c}(x_j - h_j)$
      \STATE $\log\hat{p}(x_j \mid c) \leftarrow
             \log(\hat{F}^+ - \hat{F}^- + \epsilon_s) -
             \log(2h_j + \epsilon_s)$
    \ENDIF
  \ENDFOR
  \STATE $\log\hat{p}(c \mid x) \leftarrow \log\hat{p}(c) +
         \sum_j \log\hat{p}(x_j \mid c)$
\ENDFOR
\STATE \textbf{return} $\arg\max_c \log\hat{p}(c \mid x)$
\end{algorithmic}
\end{algorithm}


% ============================================================
\section{Extended Experimental Setup}
\label{app:setup}
% ============================================================

\subsection{Hyperparameter Configurations}
\label{app:hyperparams}

Table~\ref{tab:hyperparams} lists the fixed default hyperparameters used for
every method in the main benchmark comparison. No method-specific held-out
tuning was performed; the same configuration was applied uniformly across all
datasets. For \mist{}, the defaults were chosen to reflect the theoretically
motivated operating points: $k{=}64$ (sketch capacity at which the safety
margin of the $4\varepsilon_{\mathrm{sketch}}$ bound saturates;
Appendix~\ref{app:sketch-bound}), $\alpha{=}0.6$ (well within the plateau
$\alpha \geq 0.2$ identified in Table~\ref{tab:ablation-k-alpha}),
and $\delta{=}0.10$ (consistent with the near-flat $\sqrt{\ln(1/\delta)}$
sensitivity; Table~\ref{tab:ablation-k-alpha}).

\begin{table}[h]
\caption{Fixed default configurations used in the main benchmark comparison.}
\label{tab:hyperparams}
\centering
\small
\begin{tabular}{lll}
\toprule
Method & Key hyperparameters & Policy \\
\midrule
\mist{}-G / \mist{}-K & grace=200, $\delta$=0.10, tie=0.05, $k$=64 & fixed default \\
\vfdt{} (Gini)         & grace=200, confidence=0.95, tie=0.05        & fixed default \\
EFDT / HAT (Gini)             & grace=50, split conf=$10^{-7}$, tie=0.05    & fixed default \\
Rutkowski (Gini)             & grace=200,  $\delta$=0.10, tie=0.05    & fixed default \\
\slda{} / \sqda{}      & shrinkage=$10^{-4}$                          & fixed default \\
NCM / RLS / PEC        & implementation defaults                       & fixed default \\
\bottomrule
\end{tabular}
\end{table}

\subsection{Protocol Mismatch in Drift-Reset Trees}
\label{app:noreset}

Table~\ref{tab:noreset} reports a focused diagnostic comparing EFDT and HAT
under their default drift-triggered reset configuration against the no-reset
variant used throughout the main experiments. Under strict Online CIL,
class arrivals are permanent structural events; drift detectors treat them
as transient shifts and discard the associated leaf statistics upon reset,
causing near-random performance.

When resets are disabled, HAT recovers substantial accuracy on several
datasets (e.g.\ Pendigits: $0.250 \to 0.614$, Split-MNIST: $0.170 \to 0.324$),
confirming that its failure under the default configuration is a consequence
of protocol misalignment. EFDT (no-reset) shows modest recovery on some
datasets (e.g.\ Covertype: $0.167 \to 0.277$, Pendigits: $0.196 \to 0.272$)
but remains below competitive accuracy, consistent with cold-start instability
without inheritance.
% When resets are disabled, HAT recovers substantial accuracy on all five
% datasets (e.g.\ Split-MNIST: $0.183 \to 0.324$, Synth-10: $0.333 \to 0.423$),
% confirming that its failure under the default configuration is a consequence
% of protocol misalignment rather than algorithmic failure.
% EFDT (no-reset) shows identical numbers to EFDT (default) across all five
% datasets in Table~\ref{tab:noreset}, indicating that the no-reset flag alone
% does not recover EFDT's accuracy in our implementation; we trace this to a
% baseline implementation issue whereby the EFDT-NR variant incorrectly
% inherited from \mist{}'s base class, disabling effective leaf updates.
% Corrected EFDT-NR results are reported in the main comparison
% (Table~\ref{tab:tree-results}) after fixing this issue.
% This motivates the \emph{no-reset} (NR) convention adopted for all
% tree-based baselines in the main comparison, with careful verification
% of each baseline's implementation against its original codebase.

% \begin{table}[h]
% \caption{Default vs.\ no-reset diagnostics on selected datasets (final mean
% accuracy, single seed).}
% \label{tab:noreset}
% \centering
% \small
% \begin{tabular}{lccccc}
% \toprule
% Method & Synth-10 & Iris & Wine & Covertype & Split-MNIST \\
% \midrule
% EFDT (default)   & 0.695 & 0.927 & 0.621 & 0.263 & 0.249 \\
% EFDT (no-reset)  & 1.000 & 0.933 & 0.970 & 0.550 & 0.867 \\
% HAT (default)    & 0.100 & 0.333 & 0.333 & 0.186 & 0.168 \\
% HAT (no-reset)   & 1.000 & 0.933 & 0.970 & 0.550 & 0.867 \\
% \bottomrule
% \end{tabular}
% \end{table}
\begin{table}[h]
\caption{Default vs.\ no-reset diagnostics on selected datasets (final mean
accuracy, mean over seeds).}
\label{tab:noreset}
\centering\small
\begin{tabular}{lcccccc}
\toprule
Method & Pendigits & Shuttle & Letter & HAR & Covertype & Split-MNIST \\
\midrule
EFDT+ADWIN (reset)  & 0.196 & 0.253 & 0.121 & 0.167 & 0.167 & 0.167 \\
EFDT (no-reset)     & 0.272 & 0.143 & 0.168 & 0.202 & 0.277 & 0.182 \\
HAT (default)       & 0.250 & 0.207 & 0.107 & 0.231 & 0.168 & 0.170 \\
HAT (no-reset)      & 0.614 & 0.220 & 0.237 & 0.307 & 0.262 & 0.324 \\
\bottomrule
\end{tabular}
\end{table}

\subsection{Baseline Fairness and Tuning Policy}
\label{app:baselines-fairness}

All main benchmark results (Tables~\ref{tab:tree-results}
and~\ref{tab:main-forgetting}) are produced without method-specific held-out
tuning; the policy in Table~\ref{tab:hyperparams} is applied uniformly to all
methods. For non-tree baselines, we use published implementation defaults
without further adjustment. \arf{}-NR uses 10 base trees with drift detection
disabled, matching the no-reset convention of the tree-based group.


% ============================================================
\section{Main Results: Extended Analysis}
\label{app:results}
% ============================================================

\subsection{Prequential Accuracy and Knowledge Inheritance}
\label{app:prequential}

Figure~\ref{fig:inheritance-effect} illustrates the cost of disabling
knowledge inheritance at a single representative split event (HAR dataset,
Task~3 arrival, vertical dashed line).

\textbf{Panel~(A) — cold-start on the new task.}
Immediately after the split, \mist{}-G and \mist{}-K maintain stable
accuracy on the newly arriving task (${\approx}0.75$ and ${\approx}0.88$,
respectively), because the inherited parent-node statistics provide an
informed starting distribution for the child leaf.  The \emph{w/o
Inheritance} ablation begins at the same level but degrades monotonically
to $0.00$ by the end of the window ($\Delta{=}{+}0.75$ relative to
\mist{}-G), indicating that without transferred statistics the leaf can no
longer distinguish the new task's classes once the parent's raw counts are
discarded.

\textbf{Panel~(B) — catastrophic forgetting of an old task.}
Concurrently, Task~1 accuracy (measured on a held-out test set throughout
the stream) remains at ${\approx}0.73$ for all inheritance-enabled variants
but collapses to $0.00$ for the ablation immediately after the split
($\Delta{=}{+}0.73$).  This confirms that inheritance is the mechanism
preventing catastrophic forgetting: when a leaf splits, the child nodes
must re-estimate the old-class boundary from scratch without inherited
statistics, causing complete loss of prior knowledge.

Taken together, the two panels show that a single missing inheritance step
produces a simultaneous double failure — the model cannot learn the new task
\emph{and} forgets old ones — consistent with the variance-reduction bound
in Proposition~\ref{prop:variance}.

\begin{figure}[t]
\centering
\includegraphics[width=0.9\textwidth]{mist_figures/inheritance_effect.pdf}
\caption{Effect of knowledge inheritance at the Task~3 split event on HAR.
\textbf{(A)}~New-task (Task~3) accuracy versus samples after the split:
inheritance-enabled variants (\mist{}-G, \mist{}-K) remain stable while the
\emph{w/o Inheritance} ablation degrades to $0.00$ ($\Delta{=}{+}0.75$).
\textbf{(B)}~Old-task (Task~1) accuracy over the full split window: the
ablation collapses to $0.00$ after the split while \mist{}-G retains $0.73$
accuracy ($\Delta{=}{+}0.73$).}
\label{fig:inheritance-effect}
\end{figure}
% \subsection{Prequential Accuracy and Cold-Start Mitigation}
% \label{app:prequential}

% Figure~\ref{fig:coldstart-zoom-main} zooms in on the first task-split event
% (vertical dashed line) and plots Task-0 test accuracy as a function of samples
% relative to that split. Two consistent patterns emerge across datasets.

% First, \mist{}-G with inheritance recovers faster after each class arrival
% than the \emph{w/o Inheritance} ablation, with shallower post-split accuracy
% dips whose recovery advantage scales with $n_{\mathrm{parent}}$, consistent
% with the $O(1/(\alpha n_{\mathrm{parent}}))$ variance reduction bound of
% Proposition~\ref{prop:variance}. Across datasets, inheritance reduces the
% depth of post-split accuracy dips by an average of $6.3\%$. Second,
% \mist{}-G maintains lower prediction variance than \vfdt{}: the
% McDiarmid-calibrated GapTest suppresses spurious micro-splits that
% occasionally perturb \vfdt{}'s predictions between class-arrival events.
% Full prequential curves (sliding window, $w{=}1000$) are provided in the
% supplementary materials.

% \begin{figure*}[t]
% \centering
% \hfill
% \includegraphics[width=0.49\textwidth]{mist_figures/coldstart_Covertype.pdf}
% \hfill
% \includegraphics[width=0.49\textwidth]{mist_figures/coldstart_Split-MNIST.pdf}
% \caption{Cold-start zoom at the first task split (vertical dashed line).
% Curves show Task-0 test accuracy versus samples relative to the first split.
% Inheritance keeps the post-split trajectory stable on Split-MNIST; on
% Covertype, inheritance-based variants recover faster and to higher plateaus.}
% \label{fig:coldstart-zoom-main}
% \end{figure*}

\subsection{Forgetting Results: Full Tables and Per-Task Curves}
\label{app:forgetting-full}

Table~\ref{tab:main-forgetting} reports average forgetting across all standard
benchmarks. EFDT and HAT (both NR) show elevated forgetting on Shuttle ($0.400$)
and HAR ($0.314$), consistent with Figure~\ref{fig:tree-size}: these methods
trigger more splits, creating smaller and less stable leaf predictors whose class
estimates degrade faster under new class arrivals. The negative forgetting entry
for \mist{}-G on Shuttle ($-0.035$) reflects a genuine accuracy improvement on
Task-0 as the GNB accumulator refines its per-class likelihood estimates over
the stream; no split occurs on Shuttle, so this is not attributable to structural
repartitioning.

Per-task accuracy decay curves for Synth-10, Covertype, and Split-MNIST are
provided in the supplementary materials. Across all three datasets, \mist{}-G
with inheritance retains earlier tasks substantially better than the
no-inheritance ablation: per-task accuracy curves decay more slowly and
stabilise at higher plateaus, with the gap proportional to $n_{\mathrm{parent}}$
at the triggering split event, consistent with Proposition~\ref{prop:variance}.

\begin{table*}[t]
\caption{Average forgetting (mean $\pm$ std over 5 seeds) on the standard
Online CIL benchmarks (lower is better). \textbf{Bold}: best mean within group.}
\label{tab:main-forgetting}
\centering
\resizebox{\textwidth}{!}{%
\setlength{\tabcolsep}{3pt}
\begin{tabular}{lccccccccccc}
\toprule
\textbf{Method}
  & Synth-10 & Synth-50 & Wine & Iris
  & Covertype & Split-MNIST & Pendigits & Shuttle & Letter & HAR & \textbf{Mean} \\
\midrule
\multicolumn{12}{l}{\textit{Tree-based (CIL-adapted)}} \\
\mist{}-G (ours)
  & 0.001$\pm$0.003 & 0.001$\pm$0.002 & 0.038$\pm$0.039 & 0.080$\pm$0.045
  & 0.143$\pm$0.017 & 0.051$\pm$0.003 & 0.059$\pm$0.016 & $-$0.035$\pm$0.040
  & 0.089$\pm$0.005 & 0.103$\pm$0.006 & \textbf{0.0530} \\
\mist{}-K (ours)
  & 0.001$\pm$0.001 & 0.000$\pm$0.000 & 0.055$\pm$0.045 & 0.100$\pm$0.035
  & 0.200$\pm$0.027 & 0.064$\pm$0.004 & 0.067$\pm$0.019 & 0.000$\pm$0.007
  & 0.066$\pm$0.003 & 0.161$\pm$0.007 & 0.0714 \\
\midrule
\vfdt{}
  & 0.001$\pm$0.003 & 0.750$\pm$0.001 & 0.038$\pm$0.035 & 0.080$\pm$0.040
  & 0.340$\pm$0.013 & 0.066$\pm$0.004 & 0.581$\pm$0.059 & 0.486$\pm$0.140
  & 0.612$\pm$0.072 & 0.600$\pm$0.084 & 0.3554 \\
Rutkowski
  & 0.000$\pm$0.000 & 0.000$\pm$0.000 & 0.038$\pm$0.035 & 0.080$\pm$0.040
  & 0.143$\pm$0.015 & 0.060$\pm$0.002 & 0.066$\pm$0.008 & 0.367$\pm$0.099
  & 0.089$\pm$0.004 & 0.102$\pm$0.006 & \textbf{0.0945} \\
EFDT (NR)
  & 0.024$\pm$0.050 & 0.250$\pm$0.092 & 0.038$\pm$0.035 & 0.080$\pm$0.040
  & 0.490$\pm$0.096 & 0.044$\pm$0.049 & 0.720$\pm$0.083 & 0.151$\pm$0.099
  & 0.300$\pm$0.187 & 0.804$\pm$0.122 & 0.2901 \\
HAT (NR)
  & $-$0.035$\pm$0.106 & 0.227$\pm$0.105 & 0.038$\pm$0.035 & 0.080$\pm$0.040
  & 0.740$\pm$0.128 & $-$0.038$\pm$0.105 & 0.423$\pm$0.035 & 0.443$\pm$0.113
  & 0.094$\pm$0.059 & 0.516$\pm$0.051 & 0.2488 \\
\midrule
\multicolumn{12}{l}{\textit{Non-tree baselines (reference)}} \\
\slda{}
  & 0.000$\pm$0.000 & 0.000$\pm$0.000 & 0.014$\pm$0.020 & 0.020$\pm$0.027
  & 0.202$\pm$0.009 & 0.074$\pm$0.005 & 0.084$\pm$0.009 & 0.210$\pm$0.041
  & 0.106$\pm$0.007 & 0.014$\pm$0.002 & 0.0724 \\
\sqda{}
  & 0.000$\pm$0.000 & 0.002$\pm$0.003 & 0.000$\pm$0.000 & 0.050$\pm$0.050
  & 0.226$\pm$0.002 & 0.026$\pm$0.005 & 0.011$\pm$0.004 & 0.020$\pm$0.010
  & 0.041$\pm$0.004 & 0.021$\pm$0.002 & \textbf{0.0397} \\
NCM
  & 0.000$\pm$0.000 & 0.000$\pm$0.000 & 0.007$\pm$0.016 & 0.150$\pm$0.061
  & 0.273$\pm$0.011 & 0.084$\pm$0.004 & 0.073$\pm$0.011 & 0.079$\pm$0.033
  & 0.103$\pm$0.004 & 0.143$\pm$0.007 & 0.0912 \\
RLS
  & $-$0.000$\pm$0.004 & 0.000$\pm$0.000 & 0.029$\pm$0.047 & 0.180$\pm$0.057
  & 0.219$\pm$0.007 & 0.109$\pm$0.006 & 0.066$\pm$0.012 & 0.428$\pm$0.025
  & 0.147$\pm$0.009 & 0.013$\pm$0.001 & 0.1191 \\
PEC
  & 0.001$\pm$0.001 & 0.017$\pm$0.012 & 0.152$\pm$0.067 & 0.090$\pm$0.049 & 0.231$\pm$0.007 & 0.042$\pm$0.005 & 0.026$\pm$0.003 & 0.004$\pm$0.008 & 0.086$\pm$0.004 & 0.099$\pm$0.041 & 0.0748 \\
\bottomrule
\end{tabular}}
\end{table*}

% \subsection{Extended Baseline Comparisons}
% \label{app:baselines}

% Full accuracy and forgetting results for Rutkowski, EFDT (NR), HAT (NR), and
% \arf{}-NR are included in the main benchmark tables:
% Table~\ref{tab:tree-results} (final mean accuracy) and
% Table~\ref{tab:main-forgetting} (average forgetting). We note a structural
% asymmetry between the tree and non-tree groups: non-tree methods maintain a
% single global representation across all classes, so their accuracy and
% forgetting numbers are best interpreted as a reference regime rather than a
% direct competitor. The comparison is most meaningful on near-Gaussian
% benchmarks where the underlying assumptions of SQDA and SLDA are satisfied;
% on non-Gaussian geometry the comparison reverses (Table~\ref{tab:stress}).

\subsection{Tree-Dynamics Diagnostics}
\label{app:tree-dynamics}

Figure~\ref{fig:tree-size} plots the evolution of four structural quantities
over the training stream: number of leaves, cumulative splits, weighted leaf
Gini impurity, and mean local class cardinality $K_{\mathrm{local}}$. The
comparison between the McDiarmid GapTest (default) and the Hoeffding-based
variant reveals the structural regularisation effect described in Section~4.1:
the Hoeffding criterion triggers substantially more splits as $K_{\mathrm{local}}$
grows, 
yet Table~\ref{tab:ablation-core} shows only marginal aggregate accuracy
gain ($+1.02\%$); the structurally meaningful difference lies in tree
compactness (Figure~\ref{fig:tree-size}), not raw accuracy.
The self-clearing property
of $K_{\mathrm{local}}^{(\ell)}$---visible in the lower-right panel---validates
the memory analysis of Appendix~\ref{app:memory-full}: as leaves purify,
fewer classes remain active per leaf, and sketch memory decreases accordingly.

\begin{figure}[t]
\centering
\includegraphics[width=0.48\textwidth]{mist_figures/num_leaves.pdf}
\hfill
\includegraphics[width=0.48\textwidth]{mist_figures/num_splits.pdf}\\[4pt]
\includegraphics[width=0.48\textwidth]{mist_figures/weighted_gini.pdf}
\hfill
\includegraphics[width=0.48\textwidth]{mist_figures/klocal_mean.pdf}
\caption{Tree-dynamics diagnostics on Synth-50, Covertype, and Split-MNIST
         (standard, unweighted Gini criterion for all models).
\textbf{Top:} number of leaves and cumulative splits over the stream.
\textbf{Bottom:} weighted leaf Gini impurity and mean local class cardinality
$K_{\mathrm{local}}$.
Removing the union bound over attributes (\textsc{GainTestNoUnion}) causes
systematic over-splitting---up to $6\times$ more leaves than \mist{}---without
a commensurate gain in leaf purity (bottom-left), confirming that the union
bound is essential for controlling split proliferation as $K_{\mathrm{local}}$
grows.
\textsc{VFDT} (Hoeffding bound, no inheritance) shows a comparable split rate
to \textsc{GainTest} but higher residual impurity, reflecting the absence of
knowledge inheritance.
The staircase pattern in $K_{\mathrm{local}}$ (bottom-right) marks task
boundaries; \mist{}'s tight McDiarmid radius delays new splits until sufficient
evidence accumulates across all $K_{\mathrm{local}}$ classes.}
\label{fig:tree-size}
\end{figure}


Table~\ref{tab:inheritance-ablation} compares the default truncated-Gaussian
projection (Eq.~\eqref{eq:mu_left_right}) against a
sketch-quantile projection (\textsc{SketchInherit}) that initialises child
means directly from parent KLL quantiles. Both variants produce identical
accuracy and forgetting on Split-MNIST,
confirming that the post-split cold-start benefit derives from the warm-start
mechanism itself---i.e., providing statistically grounded priors to
children---rather than from the specific formula used to project those priors.
The sketch-quantile variant is $1.27\times$ slower on average; the overhead
is not amortised over the stream since splits are infrequent. The large
per-seed variance in runtime (Gaussian: ${\pm}27.4$ s; Sketch-Quantile:
${\pm}17.5$ s) reflects differences in tree depth reached across seeds rather
than algorithmic instability. The default truncated-Gaussian projection is
therefore preferred on near-Gaussian benchmarks where both variants are
equivalent in accuracy.
\begin{table}[h]
\caption{Split-time inheritance ablation on Split-MNIST (5 seeds). The non-parametric initialisation matches
accuracy but is $1.27\times$ slower on average.}
\label{tab:inheritance-ablation}
\centering
\small
\begin{tabular}{lcccc}
\toprule
Variant & Accuracy & Forgetting & Runtime (s) & Relative time \\
\midrule
Gaussian projection (default) & $0.7717\pm0.007$ & $0.0739\pm0.005$ & $69.35\pm27.39$ & $1.00\times$ \\
Sketch-quantile projection    & $0.7717\pm0.007$ & $0.0739\pm0.005$ & $88.35\pm17.48$ & $1.27\times$ \\
\bottomrule
\end{tabular}
\end{table}

% ============================================================
\section{Ablation and Sensitivity Analysis}
\label{app:ablation}
% ============================================================

\subsection{Sensitivity to $k$, $\alpha$, and $\delta$}
\label{app:sensitivity}

Table~\ref{tab:ablation-k-alpha} reports accuracy averaged over Synth-10,
Iris, Covertype, and Split-MNIST as each hyperparameter is swept in
isolation with all others held at default. Two plateaus are clearly
visible. For sketch capacity, performance saturates at $k \geq 32$, with negligible practical variation for $k > 64$; this is consistent with the $4/k$
decay of the sketch error bound (Section~\ref{app:sketch-bound}) entering
the negligible regime. For inheritance discount, the transition is sharp:
$\alpha{=}0.0$ collapses to $0.7439$, matching the \emph{w/o Inheritance}
ablation exactly (Table~\ref{tab:ablation-core}), while any $\alpha \geq 0.2$
recovers near-full performance. For failure probability $\delta$, the
near-constant accuracy across nearly two orders of magnitude confirms the
theoretically predicted $\sqrt{\ln(1/\delta)}$ plateau in the operational
radius (Eq.~\eqref{eq:radius}): the logarithmic dependence makes
$\delta$ a low-leverage parameter in practice.

\begin{table}[h]
\caption{Sensitivity of \mist{}-G to sketch capacity $k$, inheritance
discount $\alpha$, and split-confidence level $\delta$ (average ten standard datasets). The near-identical
values across $\delta$ reflect the theoretically predicted
$\sqrt{\ln(1/\delta)}$ plateau and are a feature, not a reporting artifact
(confirmed experimentally across $\delta \in [0.01, 0.30]$).}
\label{tab:ablation-k-alpha}
\centering
\small
\begin{tabular}{lcccccc}
\toprule
$k$ & 16 & 32 & 64 & 128 & 256 & 512 \\
Avg Acc. & 0.8219 & 0.8221 & 0.8225 & 0.8225 & 0.8260 & 0.8252 \\
\midrule
$\alpha$ & 0.0 & 0.2 & 0.4 & 0.6 & 0.8 & 1.0 \\
Avg Acc. & 0.7439 & 0.8226 & 0.8226 & 0.8226 & 0.8226 & 0.8226 \\
\midrule
$\delta$ & 0.01 & 0.05 & 0.10 & 0.15 & 0.20 & 0.30 \\
Avg Acc. & 0.8226 & 0.8226 & 0.8225 & 0.8226 & 0.8226 & 0.8226 \\
\bottomrule
\end{tabular}
\end{table}


\subsection{Threshold-Density Ablation}
\label{app:threshold}
MIST generates split-threshold candidates from inter-class quantile midpoints,
yielding at most $K_{\mathrm{local}}-1$ candidates per feature (sparse policy).
Table~\ref{tab:threshold-density} evaluates denser grids of 10 and 20
candidates.
Table~\ref{tab:threshold-density} evaluates denser grids of 10 and 20
candidates. The observed non-monotone runtime (dense grids are modestly
faster on average) therefore cannot be attributed to a reduced per-candidate
cost; it more likely reflects reduced tree-growth overhead when a larger
candidate set occasionally identifies a stronger split earlier, slightly
shortening the overall training trajectory. Runtime variance increases
substantially with grid density, however (std up to 28\,s on Covertype
vs.\ ${<}2$\,s for sparse), making this benefit unreliable.
Crucially, accuracy does not change at any density level across both
datasets. Crucially, accuracy does not change at any density level across both
datasets. This confirms that (i)~the sparse quantile-midpoint policy already
identifies the same best split as a much denser grid, because KLL quantiles
concentrate around class-conditional modes; and (ii)~the union bound over
candidate thresholds in Eq.~\eqref{eq:radius} does not conservatively expand
the radius enough to suppress valid splits.
The sparse policy therefore remains preferable for its stability
(std $<$2 s vs.\ up to 28 s for dense grids on Covertype).
\begin{table}[h]
\caption{Threshold-density ablation (mean runtime; 5 seeds). Accuracy is unchanged at all densities.
Dense grids show lower mean runtime but substantially higher variance,
confirming the sparse policy as the stable, efficient default.}
\label{tab:threshold-density}
\centering
\small
\resizebox{\textwidth}{!}{%
\begin{tabular}{lccccc}
\toprule
Dataset & Sparse (1) & Dense (10) & Dense (20) & Accuracy & Acc.\ change \\
\midrule
Covertype   & $65.40\pm1.26$ s & $53.19\pm19.12$ s ($-19\%$) & $53.09\pm28.40$ s ($-19\%$) & $0.694\pm0.006$ & none \\
Split-MNIST & $126.10\pm18.00$ s & $102.24\pm9.54$ s ($-19\%$) & $99.93\pm1.63$ s ($-21\%$) & $0.772\pm0.007$ & none \\
\bottomrule
\end{tabular}%
}
\end{table}

\subsection{\mist{}-K Parameter Sensitivity}
\label{app:mistk-sensitivity}
Table~\ref{tab:mistk-sensitivity} sweeps the two \mist{}-K-specific
hyperparameters: smoothing constant $\epsilon_s$ and bandwidth multiplier
$\beta$. Across all ten benchmarks,
accuracy is nearly invariant to $\epsilon_s$: every value in
$[10^{-5}, 10^{-1}]$ yields the same average of $0.8322$.
Note that $\epsilon_s$ appears in Eq.~\eqref{eq:mistk-likelihood} for
\emph{all} features, including numerical ones; however, on purely numerical
benchmarks the sketch-CDF differences $\hat{F}^+ - \hat{F}^-$ dominate and
the additive $\epsilon_s$ contribution is negligible, making accuracy
insensitive to its value. For nominal (categorical) features, $\epsilon_s$
acts as Laplace smoothing; since all standard benchmarks are purely
numerical, this branch is inactive and the invariance to $\epsilon_s$ is
expected.
The bandwidth multiplier $\beta$ does influence accuracy: performance peaks
at $\beta \in [1.0, 2.0]$ (avg.\ acc.\ $0.8336$--$0.8338$) and
degrades at both extremes—under-smoothing at $\beta{=}0.25$ ($0.7963$,
$-4.5\%$) collapses adjacent class-conditional modes, and over-smoothing
at $\beta{=}4.0$ ($0.8141$, $-2.3\%$) blurs interclass boundaries.
We recommend $\beta{=}1.0$ as the safe default; $\epsilon_s$ is irrelevant
for numerical-only data and can be set to any value in the standard range.

\begin{table}[h]
\caption{Sensitivity of \mist{}-K average accuracy (over 10 datasets,
5 seeds) to smoothing constant $\epsilon_s$
and bandwidth multiplier $\beta$. $\epsilon_s$ has no effect because all
benchmarks are purely numerical (KLL sketch-CDF is used for likelihood;
Laplace smoothing is inactive).}
\label{tab:mistk-sensitivity}
\centering
\small
\begin{tabular}{lccccc}
\toprule
$\epsilon_s$ & $10^{-5}$ & $10^{-4}$ & $10^{-3}$ & $10^{-2}$ & $10^{-1}$ \\
Avg Acc.     & 0.8322    & 0.8322    & 0.8322    & 0.8322    & 0.8322    \\
\midrule
$\beta$  & 0.25   & 0.50   & 1.00   & 2.00   & 4.00   \\
Avg Acc. & 0.7963 & 0.8162 & 0.8336 & 0.8338 & 0.8141 \\
\bottomrule
\end{tabular}
\end{table}

% ============================================================
\section{Memory and Computational Analysis}
\label{app:memory-all}
% ============================================================

\subsection{Memory Footprint, Throughput, and Trade-offs}
\label{app:memory-full}

Let $L$ be the number of leaves, $d$ the number of features, and
$K_{\mathrm{local}}^{(\ell)}$ the number of active classes at leaf $\ell$.
With sketch capacity $k$, total sketch memory is:
\[
  M_{\mathrm{sketch}}
  = O\!\left(dk\sum_{\ell=1}^{L} K_{\mathrm{local}}^{(\ell)}\right),
\]
which is typically much smaller than the worst-case $O(dkLK)$ because
$K_{\mathrm{local}}^{(\ell)} \ll K$ once leaves become purer
(Appendix~\ref{app:memory-scaling}).

\paragraph{Transient spikes.}
Before any split occurs, the root leaf must track all classes seen so far,
inducing a transient memory spike of $O(dkK_t)$. Once the root splits, each
child receives a subset of classes and $\sum_\ell K_{\mathrm{local}}^{(\ell)}$
begins to decrease. In practice, the self-clearing property dominates within
the first $5$--$10\%$ of the stream: after initial splits, memory stabilises
below $2\times$ the final footprint. For deployments with strict peak-memory
constraints, a simple safeguard is to cap sketch creation at $K_{\mathrm{cap}}$
classes per leaf and fall back to Gaussian summaries for excess classes.

\paragraph{Throughput and Pareto frontier.}
Table~\ref{tab:memory} characterises memory and throughput on two
representative datasets that bracket the practical operating range. Among
Pareto-optimal configurations on Covertype, NCM provides the fastest and
lightest option, \slda{} provides a fast and accurate parametric baseline,
and \mist{}-K achieves the highest single-method accuracy at the cost of
greater memory and lower throughput (due to per-prediction quantile lookups
in the sketch). \mist{}-G offers a middle ground: sketch memory identical to
\mist{}-K but significantly higher throughput ($674$ vs $43$ samples/s on
Covertype), since split evaluation uses sketches but prediction uses the
faster Gaussian accumulators. EFDT (D) and HAT (D) are included for
completeness; their near-random accuracy reflects protocol misalignment
(Appendix~\ref{app:noreset}) and their throughput numbers are not meaningful
in this context.

\begin{table}[t]
\caption{Memory footprint (KB, mean over 3 seeds) and throughput (samples/s)
on two representative datasets: Covertype ($K=6$, $d=54$, low-$K$ regime)
and Split-MNIST ($K=10$, $d=50$ after PCA, higher-$K$ regime).
\textbf{Pareto-optimal} configurations on Covertype are marked with $\star$:
NCM (fastest/lightest), \slda{} (fast with high accuracy), and \mist{}-K
(highest accuracy). EFDT (D) and HAT (D) --- where (D) denotes the default
drift-reset configuration --- are included for completeness;
their near-random accuracy reflects protocol misalignment
(Appendix~\ref{app:noreset}), not resource efficiency.}
\label{tab:memory}
\centering
\small
\setlength{\tabcolsep}{3pt}
\begin{tabular}{lccccl}
\toprule
Method
  & \multicolumn{2}{c}{Memory (KB)}
  & \multicolumn{2}{c}{Throughput (samp./s)}
  & Note \\
\cmidrule(lr){2-3}\cmidrule(lr){4-5}
  & Covertype & Split-MNIST & Covertype & Split-MNIST & \\
\midrule
\mist{}-G         &  2596.6 & 12091.8 &    674 &    401 & \\
\mist{}-K $\star$ &  2596.6 & 12091.8 &     43 &     37 & highest acc \\
\vfdt{}           &   305.5 &   591.1 &   1179 &    682 & \\
\arf{}-NR         &  3044.5 &  5912.3 &   1104 &    634 & \\
EFDT (D)          &  4335.4 &      -- &     11 &     -- & near-random \\
HAT (D)           &   703.4 &      -- &    244 &     -- & near-random \\
\slda{} $\star$   &   188.5 &   243.8 &   2543 &   2123 & fast+accurate \\
\sqda{}           &   416.8 &   596.1 &   1560 &   1161 & \\
NCM $\star$       &     5.8 &     8.8 & 46778  & 31815  & ultra-fast \\
\bottomrule
\end{tabular}
\end{table}

\subsection{Self-Clearing Validation}
\label{app:memory-scaling}

To validate the self-clearing property empirically, we evaluate \mist{}-G
on multi-dimensional random-means streams ($d{=}10$, $\sigma{=}0.3$,
$K \in \{10\ldots150\}$, 2 seeds). Class means are drawn uniformly from
$[0,5]^{10}$, so splits distribute classes across multiple leaves
naturally without artificial pre-specification of the partition structure.

Table~\ref{tab:klocal} shows that $K_{\mathrm{local}}$ stabilises at
18--82\% of $K_{\mathrm{global}}$ across all tested cardinalities, consistent
with the lower-right panel of Figure~\ref{fig:tree-size} on real benchmarks.
At $K{=}150$, fewer than 30 classes remain active per leaf on average despite
the global set containing 150 classes, confirming that the self-clearing
mechanism substantially tightens the empirical memory bound relative to the
$O(dkLK)$ worst case.

\begin{table}[h]
\centering
\caption{$K_{\mathrm{local}}$ mean vs.\ $K_{\mathrm{global}}$ on
random-means streams ($d{=}10$, $\sigma{=}0.3$, 2 seeds). The
self-clearing property holds robustly across all tested class cardinalities.}
\label{tab:klocal}
\begin{tabular}{rrrr}
\toprule
$K$ & Leaves & $K_{\mathrm{local}}$ mean & Reduction \\
\midrule
25  & 20.0 &  9.0 & 64\% \\
50  & 17.5 & 13.7 & 73\% \\
75  & 33.0 & 15.5 & 79\% \\
100 & 20.5 & 29.2 & 71\% \\
150 & 21.0 & 27.1 & 82\% \\
\bottomrule
\end{tabular}
\end{table}


% ============================================================
\section{Theoretical Empirical Validation}
\label{app:theory-validation}
% ============================================================

\subsection{Empirical Validation of the $4\varepsilon_{\mathrm{sketch}}$ Bound}
\label{app:sketch-bound}

We measured $|\hat{F} - F|$ (sketched vs.\ exact Gini gain) at every split
decision under the strict Online CIL protocol across 7 datasets, 4 sketch
capacities ($k \in \{16, 32, 64, 128\}$), and 5 seeds. All results use the
corrected CIL class-incremental presentation and standard (unweighted) Gini
gain, consistent with Theorem~\ref{thm:tight}.


\paragraph{Per-dataset results at $k{=}64$.}

Table~\ref{tab:sketch-bound-k64} reports empirical sketch error across all
validated datasets at $k{=}64$, giving $4\varepsilon = 4/64 = 0.0625$.
The bound holds at every recorded split event. Safety margins range from
$2.3\times$ (Pendigits, the most marginal case) to $27.9\times$ (HAR, where
class-conditional distributions are approximately isotropic and the sketch
is highly accurate). Covertype exceeds the KLL rank-error threshold
($\varepsilon_{\mathrm{sketch}} = 1/k$) as a CDF \emph{value} error, but
this does not violate the bound: Theorem~\ref{thm:composed} controls Gini
gain error via rank error, not CDF value error directly, and
Gini err max $= 0.015 \ll 0.0625$ confirms the bound holds.

\begin{table}[h]
\centering
\caption{Empirical sketch error at $k{=}64$ ($4\varepsilon = 0.0625$).
         \emph{Events}: split decisions recorded (5 seeds, $\times$5).
         \emph{Ratio}: $4\varepsilon \,/\,$ Gini err max (worst-case
         safety margin).
         Letter is excluded; see note below.}
\label{tab:sketch-bound-k64}
\resizebox{\linewidth}{!}{%
\begin{tabular}{lrrrrrrr}
\toprule
Dataset & Events
        & CDF err mean & CDF err max
        & Gini err mean & Gini err max
        & $4\varepsilon$ & Ratio \\
\midrule
Pareto               & 20 & 0.0031 & 0.0067 & 0.00498 & 0.01126 & 0.0625 & $5.5\times$  \\
LogNormal            & 20 & 0.0080 & 0.0145 & 0.01056 & 0.02955 & 0.0625 & $2.1\times$  \\
Exponential          & 20 & 0.0047 & 0.0101 & 0.00523 & 0.01226 & 0.0625 & $5.1\times$  \\
Outlier 5\%          & 20 & 0.0068 & 0.0122 & 0.00914 & 0.02318 & 0.0625 & $2.7\times$  \\
Pendigits            &  4 & 0.0073 & 0.0145 & 0.01331 & 0.02715 & 0.0625 & $2.3\times$  \\
HAR                  &  5 & 0.0009 & 0.0014 & 0.00114 & 0.00224 & 0.0625 & $27.9\times$ \\
Shuttle              & -- & 0.0001 & 0.0001 & --      & --      & 0.0625 & --           \\
Covertype$^\ddagger$ & 12 & 0.0248 & 0.0585 & 0.00487 & 0.01486 & 0.0625 & $4.2\times$  \\
\bottomrule
\end{tabular}}
\smallskip
\begin{minipage}{\linewidth}
\small
$^\ddagger$\,CDF err max $= 0.059$ on Covertype exceeds $1/k = 0.0156$
as a CDF \emph{value} error, but this does not violate the theoretical
bound: $4\varepsilon_{\mathrm{sketch}}$ in Theorem~\ref{thm:composed}
bounds Gini gain error via \emph{rank} error, not CDF value error.
Gini err max $= 0.015 \ll 0.0625$ confirms the bound holds.
\end{minipage}
\end{table}

\paragraph{Note on zero-split datasets.}
Letter produced no split events at any $k$ under the default
\texttt{grace\_period}$=200$, consistent with the main experiment
(Table~\ref{tab:sketch-vs-gaussian-real}). 
Shuttle produced no split events under the unweighted-Gini validation
protocol: its heavy class imbalance (${\approx}80\%$ majority class)
causes the Gini gain gap to fall below the McDiarmid radius, so the
bound is vacuously satisfied.
% Shuttle likewise produced
% no split events under the validation protocol and is therefore omitted
% from the Gini error analysis.

\paragraph{Scaling with sketch capacity $k$.}

Table~\ref{tab:sketch-bound-kscale} reports safety margins across sketch
capacities for the four synthetic validation datasets, where sufficient split
events (20 per configuration) are available at all $k$ values.
The bound $4\varepsilon_{\mathrm{sketch}} \leq 4/k$ holds at every recorded
split event across all configurations. Margins degrade gracefully as $k$
decreases, consistent with the $O(1/k)$ rank-error guarantee of KLL: even at $k{=}16$ (the smallest tested), the minimum margin remains above
$2\times$. At $k{=}64$, worst-case margins range from $2.1\times$
(LogNormal) to $5.5\times$ (Pareto), confirming the bound holds with
meaningful slack across all tested configurations.

\begin{table}[h]
\centering
\caption{Ratio $4\varepsilon \,/\,$ Gini err max across $k$ values,
synthetic datasets only (20 split events each).
The $4\varepsilon_{\mathrm{sketch}}$ bound holds at every recorded event;
safety margin degrades gracefully as $k$ decreases but remains $\geq 2.1\times$
at $k{=}64$.}
\label{tab:sketch-bound-kscale}
\begin{tabular}{rrrrrrr}
\toprule
$k$ & $4\varepsilon$ & Pareto & LogNormal & Exponential & Outlier 5\% \\
\midrule
16  & 0.250  & $14\times$ & $20\times$ & $20\times$ & $9\times$  \\
32  & 0.125  & $14\times$ & $12\times$ & $25\times$ & $14\times$ \\
64  & 0.0625 & $5.5\times$ & $2.1\times$ & $5.1\times$ & $2.7\times$ \\
% 64  & 0.0625 & $13\times$ & $6\times$  & $12\times$ & $7\times$  \\
128 & 0.0312 & $13\times$ & $8\times$  & $16\times$ & $10\times$ \\
\bottomrule
\end{tabular}
\end{table}

\subsection{Sketch-Quantile vs.\ Truncated-Gaussian Inheritance on Non-Gaussian Streams}
\label{app:sketch-vs-gaussian}

This appendix evaluates the robustness of the truncated-Gaussian projection
used in Bayesian inheritance (Section~4.2) when the underlying class-conditional
distributions violate the Gaussianity assumption. We compare \mist{}-G against
\textsc{SketchInherit}, a variant that initialises child means directly from
parent KLL quantiles: the left-child mean for class $c$ on the split feature
is set to $\widehat{F}_{j,c}^{-1}(r_L/2)$, and the right-child mean to
$\widehat{F}_{j,c}^{-1}(r_L + r_R/2)$, where $r_L$ and $r_R$ are the
estimated mass fractions routing to each child. Quantiles are resolved by
nearest-point lookup with boundary clipping to $[0.01, 0.99]$. Non-split
feature parameters and class mass are inherited identically to \mist{}-G.
All other components are identical, including $\alpha{=}0.6$. Results
are averaged over 5 seeds $\{42, 123, 456, 789, 3330\}$ with
\texttt{grace\_period}$=200$, $\delta=0.10$, $k=64$.

\paragraph{Synthetic non-Gaussian streams.}
Four distributions isolate specific failure modes of the truncated-Gaussian
projection. \textbf{Pareto}: power-law tails ($\alpha{=}1.5$) with
\texttt{log1p} normalisation; tests whether log-transform approximately
Gaussianises the truncated region. \textbf{LogNormal}: log-normal marginals
($\sigma \in [0.8, 1.5]$); right-skewed distributions where
$\mathbb{E}[X \mid X \leq t]$ lies below the Gaussian prediction.
\textbf{Outlier 5\%}: Gaussian base with $5\%$ extreme outliers
(${\times}50\sigma$); tests robustness of Welford variance estimates under
contamination. \textbf{Exponential}: exponential marginals
($\lambda \in [0.5, 2]$); moderately right-skewed, between LogNormal and
Gaussian in severity. Each dataset contains $K{=}6$ classes, $d{=}10$
features, $n{=}12{,}000$ training samples ($2{,}000$ per class), presented
in strict Online CIL order.

Table~\ref{tab:sketch-vs-gaussian-synthetic} shows that gains are largest
where the Gaussian truncation assumption fails most severely: $+33.5$ pp on
Outlier~5\% (a single outlier corrupts the Welford variance estimate, making
the moment projection unreliable, while the sketch CDF estimate remains
robust) and $+2.3$ pp on LogNormal. The Pareto case ($-1.2$ pp) is the sole
exception: \texttt{log1p} normalisation approximately symmetrises the
truncated region, making the Gaussian approximation adequate while the
sketch introduces marginal quantile-grid variance.

\begin{table}[h]
\centering
\caption{Accuracy comparison on synthetic non-Gaussian streams
         (mean $\pm$ std over 5 seeds).
         $\Delta\mathrm{acc} = $ \textsc{SketchInherit} $-$ \mist{}-G.}
\label{tab:sketch-vs-gaussian-synthetic}
\resizebox{\linewidth}{!}{%
\begin{tabular}{lllrrr}
\toprule
Dataset & Distribution & \mist{}-G & \textsc{SketchInherit} & $\Delta$acc & Splits \\
\midrule
Pareto      & Power-law ($\alpha{=}1.5$), \texttt{log1p} transform
            & $0.9626\pm0.003$ & $0.9503\pm0.004$ & $-0.012$ & 5.0 \\
LogNormal   & Log-normal ($\sigma \in [0.8, 1.5]$)
            & $0.8049\pm0.017$ & $0.8277\pm0.011$ & $+0.023$ & 5.0 \\
Outlier 5\% & Gaussian $+5\%$ outliers (${\times}50\sigma$)
            & $0.3577\pm0.026$ & $0.6928\pm0.039$ & $\mathbf{+0.335}$ & 5.0 \\
Exponential & Exponential ($\lambda \in [0.5, 2]$)
            & $0.8676\pm0.025$ & $0.8809\pm0.031$ & $+0.013$ & 4.4 \\
\bottomrule
\end{tabular}}
\end{table}

\paragraph{Real-world benchmark datasets.}
Table~\ref{tab:sketch-vs-gaussian-real} confirms that on near-Gaussian data
(Pendigits, Split-MNIST, Covertype, HAR), both variants produce identical
accuracy: the truncated-Gaussian projection is harmless when valid. 
Shuttle produced no split events under either variant (Table~\ref{tab:sketch-vs-gaussian-real}),
so the bound is vacuously satisfied and no accuracy comparison is possible.
Letter
produces no splits under the default settings and is included for completeness.

\textsc{SketchInherit} is $1.27\times$ slower on average
(Table~\ref{tab:inheritance-ablation}, Relative time column),
confirming that the cold-start benefit derives from the warm-start mechanism
rather than the projection formula. The default \mist{}-G uses
truncated-Gaussian projection on near-Gaussian benchmarks where it is
equivalent at lower computational cost.

\begin{table}[h]
\centering
\caption{Accuracy comparison on standard Online CIL benchmarks
         (mean $\pm$ std over 5 seeds).
         $K$: number of classes; $d$: feature dimension.}
\label{tab:sketch-vs-gaussian-real}
\begin{tabular}{llrrrr r}
\toprule
Dataset & $K$ & $d$ & \mist{}-G & \textsc{SketchInherit} & $\Delta$acc & Splits \\
\midrule
Pendigits   & 10 & 16  & $0.8886\pm0.023$ & $0.8886\pm0.023$ & $0.000$  & 1.4 \\
Shuttle     & 7  & 9   & $0.7331\pm0.033$ & $0.7331\pm0.054$ & $0.000$  & 0.0 \\
Letter      & 26 & 16  & $0.6775\pm0.006$ & $0.6775\pm0.006$ & $0.000$  & 0.0 \\
HAR         & 6  & 561 & $0.6074\pm0.013$ & $0.6074\pm0.013$ & $0.000$  & 1.0 \\
Covertype   & 7  & 54  & $0.5590\pm0.013$ & $0.5590\pm0.013$ & $0.000$  & 0.0 \\
Split-MNIST & 10 & 50  & $0.8629\pm0.003$ & $0.8629\pm0.002$ & $0.000$  & 1.0 \\
\bottomrule
\end{tabular}
\end{table}


% ============================================================
\section{Stress-Test Dataset Descriptions}
\label{app:stress}
% ============================================================

All stress-test datasets are evaluated under the same Online CIL protocol
as the main experiments (Section~\ref{sec:setup}) with a fixed random seed; generation
code is provided in the supplementary materials.
Figure~\ref{fig:stress-vis} shows 2D PCA projections of all eight
distributions.

\paragraph{Multimodal ($K{=}8$, $d{=}10$, 4 tasks).}
Each class is drawn from a mixture of $M{=}3$ well-separated Gaussian modes.
Tests whether leaf predictors can model non-unimodal class-conditional
distributions. \mist{}-K's sketch-based marginal estimator captures the
multimodal structure without raw-data retention, while GNB-leaf methods
incur systematic bias.

\paragraph{Skewed ($K{=}8$, $d{=}10$, 4 tasks).}
Each class has a dominant mode (90\%) and a spatially separated minority
mode (10\%). The skew places the class-conditional mean outside the dominant
mode's mass, exposing the truncated-Gaussian inheritance assumption to bias.

\paragraph{AngularSectors ($K{=}4$, $d{=}8$, 2 tasks).}
Samples lie on a 2D ring; class identity is determined by angular sector.
All $K$ classes share mean $\approx\mathbf{0}$ and identical covariance, so
QDA, LDA, and Naive Bayes degrade to chance ($1/K$). A tree with
axis-aligned splits on the two signal features can progressively narrow
sectors. \mist{}-K is the only tested configuration that rises substantially
above this near-chance regime.

\paragraph{Antipodal ($K{=}8$, $d{=}10$, 4 tasks).}
Each class consists of two diametrically opposite clusters on a
hypersphere. \slda{} fails ($0.171$) because a shared covariance hyperplane
cannot separate antipodal classes; \sqda{} succeeds because per-class
Gaussian means differ.

\paragraph{Ring ($K{=}6$, $d{=}8$, 3 tasks).}
Classes are distributed on concentric hypersphere shells. The radial
direction is near-Gaussian, so most methods perform competitively. This
confirms that \mist{}'s advantage is regime-dependent: where the Gaussian
assumption is adequate, global parametric methods remain competitive.

\paragraph{HeavyTail ($K{=}8$, $d{=}10$, 4 tasks).}
Class-conditional distributions follow Student-$t$ ($\nu{=}2$) noise around
random centers. Fat tails inflate Welford variance estimates, biasing
\mist{}-G's truncated-Gaussian projection; \mist{}-K mitigates this via the
empirical sketch-based CDF.

\paragraph{NoisyFeatures ($K{=}8$, $d{=}20$, 4 tasks).}
Only $d_{\mathrm{info}}{=}2$ dimensions carry signal; the remaining 18 are
independent noise. Tests split-selection robustness under high irrelevant-
feature contamination. The conservative McDiarmid GapTest is particularly
relevant: uninformative features produce near-zero Gini gain, and a tight
confidence radius suppresses spurious splits on noise features.

\paragraph{ConceptDrift ($K{=}8$, $d{=}10$, 4 tasks).}
Class-conditional means shift linearly across tasks (drift $= 1.5$ per task
along feature~0). Included to verify that \mist{}'s conservative split
criterion does not over-commit to stale boundaries. High accuracy across all
methods ($> 0.988$) confirms that gradual drift is absorbed by the existing
split mechanism without explicit drift detection.

\begin{figure}[t]
  \centering
  \begin{subfigure}[t]{0.24\textwidth}
    \includegraphics[width=\textwidth]{mist_figures/pca/Multimodal_pca3d.png}
    \caption{Multimodal}
  \end{subfigure}
  \hfill
  \begin{subfigure}[t]{0.24\textwidth}
    \includegraphics[width=\textwidth]{mist_figures/pca/Skewed_pca3d.png}
    \caption{Skewed}
  \end{subfigure}
  \hfill
  \begin{subfigure}[t]{0.24\textwidth}
    \includegraphics[width=\textwidth]{mist_figures/pca/AngularSectors_pca3d.png}
    \caption{AngularSectors}
  \end{subfigure}
  \hfill
  \begin{subfigure}[t]{0.24\textwidth}
    \includegraphics[width=\textwidth]{mist_figures/pca/Antipodal_pca3d.png}
    \caption{Antipodal}
  \end{subfigure}
  \\[6pt]
  \begin{subfigure}[t]{0.24\textwidth}
    \includegraphics[width=\textwidth]{mist_figures/pca/Ring_pca3d.png}
    \caption{Ring}
  \end{subfigure}
  \hfill
  \begin{subfigure}[t]{0.24\textwidth}
    \includegraphics[width=\textwidth]{mist_figures/pca/HeavyTail_pca3d.png}
    \caption{HeavyTail}
  \end{subfigure}
  \hfill
  \begin{subfigure}[t]{0.24\textwidth}
    \includegraphics[width=\textwidth]{mist_figures/pca/NoisyFeatures_pca3d.png}
    \caption{NoisyFeature}
  \end{subfigure}
  \hfill
  \begin{subfigure}[t]{0.24\textwidth}
    \includegraphics[width=\textwidth]{mist_figures/pca/ConceptDrift_pca3d.png}
    \caption{ConceptDrift}
  \end{subfigure}
  \caption{2D PCA visualisations of the eight stress-test streams (50 samples
  per class shown). Each panel illustrates the geometric property that
  motivates the dataset: non-unimodal class structure (Multimodal, Skewed),
  manifold geometry (AngularSectors, Ring), cluster arrangement (Antipodal),
  and distributional properties (HeavyTail, NoisyFeature, ConceptDrift).}
  \label{fig:stress-vis}
\end{figure}


% ============================================================
\section{Mathematical Proofs}
\label{app:proofs}
% ============================================================

Table~\ref{tab:assumptions} summarises the assumptions underlying each
theoretical result. The key structural observation is that McDiarmid's
inequality requires only bounded differences, not i.i.d.; samples at a
leaf need not be independent across time, only the sensitivity constant
$c_i = 4/n$ must hold per replacement. This makes the analysis applicable
under the per-leaf i.i.d.\ protocol of~\cite{DomingosHulten2000} without requiring
global stream independence.

\begin{table}[h]
\centering
\caption{Assumptions per theoretical result.}
\label{tab:assumptions}
\begin{tabular}{llp{5cm}}
\toprule
Result & Assumption & Note \\
\midrule
Theorem~\ref{thm:tight}
  & Bounded differences per leaf
  & McDiarmid requires only bounded differences, not i.i.d.; samples at a
    leaf need not be independent across time, only $c_i = 4/n$ must hold per
    replacement \\
Theorem~\ref{thm:composed}
  & Bounded differences + KLL rank guarantee
  & Statistical term uses $\mathbb{E}[F]$ over the empirical leaf distribution;
    independence within a leaf's $n$ samples is assumed for the McDiarmid
    application, consistent with the per-leaf i.i.d.\ protocol
    of~\cite{DomingosHulten2000} \\
Theorem~\ref{thm:composed}
  & KLL rank guarantee
  & Holds w.p.\ $\geq 1-\delta_{\mathrm{sk}}$ \\
Corollary~\ref{cor:opt-displacement}
  & Gaussian class-conditionals
  & \mist{}-K replaces $\sigma$ with sketch IQR \\
Proposition~\ref{prop:variance}
  & Dirichlet-multinomial; $\alpha n_c^{\mathrm{parent}} > 2$
  & Standard conjugate models; condition ensured by conservative splits \\
\bottomrule
\end{tabular}
\end{table}

\subsection{Truncated-Gaussian Moment Updates}
\label{app:truncgauss}

When $X \sim \mathcal{N}(\mu, \sigma^2)$ is conditioned on $X \leq v$
(left child), let $\zeta = (v-\mu)/\sigma$. The exact conditional
moments are:
\begin{align}
  \mathbb{E}[X \mid X \leq v]
    &= \mu - \sigma\frac{\phi(\zeta)}{\Phi(\zeta)}, \label{eq:trunc-mean-left}\\
  \mathrm{Var}(X \mid X \leq v)
    &= \sigma^2\!\left[1 - \frac{\zeta\phi(\zeta)}{\Phi(\zeta)}
      - \left(\frac{\phi(\zeta)}{\Phi(\zeta)}\right)^{\!2}\right].
      \label{eq:trunc-var-left}
\end{align}

For the right child ($X \geq v$), replacing $\Phi(\zeta)$ with
$1{-}\Phi(\zeta)$ throughout gives:
\begin{align}
  \mathbb{E}[X \mid X \geq v]
    &= \mu + \sigma\frac{\phi(\zeta)}{1-\Phi(\zeta)}, \label{eq:trunc-mean-right}\\
  \mathrm{Var}(X \mid X \geq v)
    &= \sigma^2\!\left[1 \mathbf{+} \frac{\zeta\phi(\zeta)}{1-\Phi(\zeta)}
      - \left(\frac{\phi(\zeta)}{1-\Phi(\zeta)}\right)^{\!2}\right].
      \label{eq:trunc-var-right}
\end{align}

\noindent\textbf{Critical sign difference.}
The linear term $\zeta\phi(\zeta)/\tilde\Phi$ carries a \emph{negative}
sign for the left child (Eq.~\eqref{eq:trunc-var-left}) and a
\emph{positive} sign for the right child (Eq.~\eqref{eq:trunc-var-right}).
Algorithm~\ref{alg:split} Line~7--8 applies this conditional sign.
Using the same sign for both children would yield the incorrect formula
$[1 - \zeta\phi(\zeta)/(1{-}\Phi(\zeta)) - \cdots]$ for the right child,
which can produce negative variances (e.g.\ at $\zeta=0$, the bracket
evaluates to $1 - 0 - 1/(2\pi) \approx 0.84 > 0$ for the left but to
$1 - 0 - 1/(2\pi)$ for the wrong-sign right formula, which happens to be
positive at $\zeta=0$ but becomes negative for $\zeta \gg 0$).

Both conditional variances satisfy
$\mathrm{Var}(X \mid X \gtrless v) \leq \sigma^2$,
confirming that the inheritance projection cannot inflate
feature-parameter uncertainty. For categorical features, no moment
projection is applied; see Section~\ref{sec:inheritance} for the
empirical co-occurrence formula.

\subsection{Proof of Theorem~\ref{thm:tight} (Tightness of the Gini Sensitivity)}
\label{app:proof-tight}

We first state two supporting lemmas and then give the three-step proof.

\begin{lemma}[Gini gain is Lipschitz in threshold]
\label{lem:gini-lipschitz}
For a fixed sample set $S$, the Gini split gain $F(S;\,v)$ as a function
of the split threshold $v$ is Lipschitz with constant
$L_v = O(1/\sigma_{\min})$, where $\sigma_{\min}$ is the minimum
class-conditional standard deviation at the leaf.
\end{lemma}

\begin{proof}
A threshold displacement $\Delta v$ reroutes at most
$\Delta_{\mathrm{total}} = n\hat{p}(v)|\Delta v|$ samples in total,
where $\hat{p}(v)$ is the empirical density at $v$.
By Lemma~\ref{lem:routing-lipschitz},
$|F(S') - F(S)| \leq (4/n)\cdot\Delta_{\mathrm{total}} = 4\hat{p}(v)|\Delta v|$.
Under a Gaussian class-conditional model,
$\hat{p}(v) \leq 1/(\sigma_{\min}\sqrt{2\pi})$,
giving $L_v = 4\hat{p}(v) = O(1/\sigma_{\min})$.
\end{proof}

\begin{lemma}[Gini gain Lipschitz in routing counts]
\label{lem:routing-lipschitz}
For any perturbation of routing counts of total magnitude
$\Delta_{\mathrm{total}} = \sum_c|\Delta n_{c,L}|$, the Gini split gain
satisfies $|F(S') - F(S)| \leq (4/n)\cdot\Delta_{\mathrm{total}}$.
\end{lemma}
\begin{proof}
By Theorem~\ref{thm:tight}, a single-sample replacement (which changes
routing counts by at most 1 in total) alters $F$ by at most $4/n$.
For a total perturbation of magnitude $\Delta_{\mathrm{total}}$, applying
the triangle inequality over at most $\Delta_{\mathrm{total}}$ unit steps
gives $|F(S') - F(S)| \leq (4/n)\cdot\Delta_{\mathrm{total}}$.
\end{proof}

\textbf{Step 1: Sensitivity of node impurity.}
For any sample set $U$ of size $m$, replacing one sample changes at most
two class proportions, each by $1/m$. Since $G(U) = 1 - \sum_c p_c^2$,
and at most two terms change with $p_c + p_c' \leq 2$, we have
$|G(U) - G(U^{(i)})| \leq 2/m$.

\begin{lemma}[Node-impurity sensitivity is tight]
\label{lem:impurity-tight}
The bound $|G(U) - G(U^{(i)})| \leq 2/m$ is tight: for any $\gamma > 0$,
there exists $U$ of size $m$ and a replacement such that
$|G(U) - G(U^{(i)})| > 2/m - \gamma$.
\end{lemma}
\begin{proof}
Consider $K = 2$ with $U$ containing $n_0 = 1$ sample of class~0 and
$n_1 = m-1$ samples of class~1. Replace the single class-0 sample with
class~1. Direct computation gives:
\[
  |G(U) - G(U^{(i)})| = \frac{2(m-1)}{m^2}
  = \frac{2}{m} - \frac{2}{m^2} \to \frac{2}{m} \quad \text{as } m \to \infty.
\]
For any $\gamma > 0$, choose $m > 2/\gamma$.
\end{proof}

\textbf{Step 2: Case decomposition.}
\textit{Case~A (same routing):} both $z_i$ and $z_i'$ route to child $L$.
Then $\Delta F = \Delta G_P - (n_L/n)\Delta G_L$, and by Step~1,
$|\Delta F| \leq 2/n + (n_L/n)(2/n_L) = 4/n$.
\textit{Case~B (cross routing):} the replacement changes routing between
children. Direct expansion gives $n|\Delta F| \leq 2 < 4$
(full derivation in Appendix~\ref{app:case-b}).
Hence $|F(S) - F(S^{(i)})| \leq 4/n$ in all cases.

\textbf{Step 3: Tightness.}
Consider binary labels with $m = \lfloor\sqrt{n/2}\rfloor$ left-child
samples (all class~1) and right-child samples (all class~0); flip one
left sample to class~0. The scaled deviation
$H(m,n) = 2(n-m)(2m-1)/(nm) \to 4$ as $n \to \infty$.
For any $\gamma > 0$, choosing $n$ large enough yields
$|F(S) - F(S^{(i)})| > 4/n - \gamma$.

Steps~1--2 establish $c_i \leq 4/n$; Step~3 confirms $4/n$ is sharp
via Lemma~\ref{lem:impurity-tight}, completing the proof. \hfill$\square$

\subsection{Proof of Theorem~\ref{thm:composed} (Composed Split Error)}
\label{app:proof-composed}

By the triangle inequality,
$|\widehat{F} - F^*| \leq |F(S) - F^*| + |\widehat{F} - F(S)|$.

\textit{Statistical term.}
From Theorem~\ref{thm:tight}, $c_i = 4/n$ for all $i$.
Applying McDiarmid's inequality with a union bound over $d$ features
and $m$ thresholds:
\[
  \Pr\!\left[\max_{j,v}|F_{j,v}(S) - F^*_{j,v}| \geq \eta\right]
  \leq 2dm\exp\!\left(-\frac{n\eta^2}{8}\right).
\]
Setting the right-hand side to $\delta$ and solving gives
$\eta = \sqrt{8\ln(2dm/\delta)/n}$.

\textit{Sketch term.}
On the KLL success event (probability $\geq 1-\delta_{\mathrm{sketch}}$),
rank error $\varepsilon_{\mathrm{sketch}}$ implies
$|\hat{n}_{c,L} - n_{c,L}| \leq n_c\varepsilon_{\mathrm{sketch}}$, so
$\sum_c|\hat{n}_{c,L} - n_{c,L}| \leq n\varepsilon_{\mathrm{sketch}}$.
By Lemma~\ref{lem:routing-lipschitz},
$|\widehat{F} - F(S)| \leq (4/n)\cdot n\varepsilon_{\mathrm{sketch}}
= 4\varepsilon_{\mathrm{sketch}}$.

Combining via union bound gives confidence $\geq 1-\delta-\delta_{\mathrm{sketch}}$.
\hfill$\square$

\subsection{Proof of Corollary~\ref{cor:opt-displacement}
           (Optimization Displacement)}
\label{app:proof-displacement}

The displacement bound follows a three-term decomposition.
Throughout, we work under the KLL success event (probability
$\geq 1-\delta_{\mathrm{sketch}}$) on which the rank-error guarantee
$\varepsilon_{\mathrm{sketch}}$ holds.

\textbf{Step~(i): Quantile displacement.}
Under Gaussianity of $X_j \mid c$, a KLL rank error of
$\varepsilon_{\mathrm{sketch}}$ implies that the empirical median
$\hat{q}_{0.5,c}$ differs from the true median $q_{0.5,c}$ by at most
$\varepsilon_{\mathrm{sketch}}$ in probability mass.
By the Mean Value Theorem applied to the class-conditional CDF
$\Phi_c$, the spatial displacement satisfies
\[
  |\hat{q}_{0.5,c} - q_{0.5,c}|
  = \frac{\varepsilon_{\mathrm{sketch}}}{p_c(\tilde{q})},
\]
for some $\tilde{q}$ between $q_{0.5,c}$ and $\hat{q}_{0.5,c}$,
where $p_c$ is the class-conditional density.
An upper bound on this displacement therefore requires a
\emph{lower} bound on $p_c(\tilde{q})$ over the displaced interval.

\textit{Small-displacement regime.}
We work in the regime $|\hat{q}_{0.5,c}-q_{0.5,c}| \ll \sigma_c$,
i.e.\ the displacement is small relative to the class spread.
In this regime the density $p_c$ is approximately constant over the
interval $[\min(q,\hat{q}),\max(q,\hat{q})]$ at its value at the
median, which for a Gaussian satisfies
$p_c(q_{0.5,c}) = 1/(\sigma_c\sqrt{2\pi})$.
Hence $p_{\min} \approx 1/(\sigma_c\sqrt{2\pi})$, giving
\[
  |\hat{q}_{0.5,c} - q_{0.5,c}|
  \;\lesssim\;
  \varepsilon_{\mathrm{sketch}} \cdot \sigma_c\sqrt{2\pi}.
\]
This bound is an asymptotic statement as
$\varepsilon_{\mathrm{sketch}}\to 0$: for small rank errors the
approximation is tight; for large $|\zeta|$ where the tail density
drops significantly below the mode density, the bound becomes
conservative.
The practical impact is negligible, as confirmed by the
$\leq 4{\times}$ empirical safety margins in
Table~\ref{tab:sketch-bound-kscale}~\citep{Karnin2016}.

\textbf{Step~(ii): Threshold displacement.}
Taking $\sigma_{\max} = \max_c \sigma_c$ and the result of Step~(i):
\[
  |\hat{v}^* - v^*|
  \;\leq\;
  \varepsilon_{\mathrm{sketch}} \cdot \sigma_{\max}\sqrt{2\pi}.
\]

\textbf{Step~(iii): Gini loss from displacement.}
By Lemma~\ref{lem:gini-lipschitz}, the empirical Gini split gain
$F(S;\cdot)$ is Lipschitz in the threshold with constant
$L_v = O(1/\sigma_{\min})$.
With condition number $\kappa = \sigma_{\max}/\sigma_{\min}$:
\[
  F(S;\,v^*) - F(S;\,\hat{v}^*)
  \;\leq\; L_v \cdot |\hat{v}^* - v^*|
  \;\leq\; 4\kappa\,\varepsilon_{\mathrm{sketch}}.
\]
This bounds the \emph{empirical} Gini loss incurred by substituting
the sketch-derived threshold $\hat{v}^*$ for the exact optimum $v^*$
on the same sample set $S$.

\textbf{Combining into Eq.~\eqref{eq:composed-full}.}
Equation~\eqref{eq:composed-full} controls the total deviation of the
sketch-evaluated gain at $\hat{v}^*$ from the population-optimal gain
$F^*(v^*)$.
By the triangle inequality:
\[
  \bigl|F(S;\,\hat{v}^*) - F^*(v^*)\bigr|
  \;\leq\;
  \underbrace{\bigl|F(S;\,v^*) - F^*(v^*)\bigr|}_{\text{statistical}}
  +\;
  \underbrace{\bigl[F(S;\,v^*) - F(S;\,\hat{v}^*)\bigr]}_{\text{displacement}}.
\]
The statistical term is bounded by
$\sqrt{8\ln(2dm/\delta)/n} + 4\varepsilon_{\mathrm{sketch}}$
via Theorem~\ref{thm:composed} (holding with probability
$\geq 1-\delta-\delta_{\mathrm{sketch}}$).
The displacement term is bounded by
$4\kappa\,\varepsilon_{\mathrm{sketch}}$ from Step~(iii).
Summing gives:
\[
  \bigl|F(S;\,\hat{v}^*) - F^*(v^*)\bigr|
  \;\leq\;
  \sqrt{\frac{8\ln(2dm/\delta)}{n}}
  + 4(1+\kappa)\,\varepsilon_{\mathrm{sketch}},
\]
which is Eq.~\eqref{eq:composed-full}.
In particular, since $F(S;v^*)-F(S;\hat{v}^*)\leq 4\kappa\varepsilon_{\mathrm{sketch}}$
(Step~iii) and $4\kappa\varepsilon_{\mathrm{sketch}} \leq$ the RHS of
Eq.~\eqref{eq:composed-full}, the stated bound also covers the
empirical displacement loss directly.
\hfill$\square$

\subsection{Proof of Proposition~\ref{prop:variance}
           (Post-Split Variance Reduction)}
\label{app:proof-variance}

\textit{Part~(i): Class probabilities.}
The inherited prior is Dirichlet with total mass
$S_0 = \alpha n_{\mathrm{parent}}$.
After $t$ local samples, $S_t = S_0 + t$.
The Dirichlet posterior variance:
\[
  \mathrm{Var}(\theta_c \mid \mathcal{D}_t)
  = \frac{\beta_c^{(t)}(S_t - \beta_c^{(t)})}{S_t^2(S_t+1)}
  \leq \frac{1}{4(S_t+1)}
  \leq \frac{1}{4(S_0+1)},
\]
where the first inequality uses the AM-GM bound
$\beta_c^{(t)}(S_t-\beta_c^{(t)}) \leq S_t^2/4$.

\textit{Part~(ii): Feature parameters.}
Under the NIG conjugate model with
$n_0 = \alpha n_c^{\mathrm{parent}}$ pseudo-observations,
the marginal posterior on $\mu_j^c$ at child creation is a
Student-$t$ distribution with $\nu = n_0$ degrees of freedom and
scale $(\sigma_j^c)^2/n_0$.
Its variance exists when $n_0 > 2$ and equals:
\[
  \frac{\nu}{\nu-2}\cdot\frac{(\sigma_j^c)^2}{n_0}
  = \frac{(\sigma_j^c)^2}{n_0-2}
  = \frac{(\sigma_j^c)^2}{\alpha n_c^{\mathrm{parent}} - 2}.
\]
For the split feature, substituting
$\sigma_{\mathrm{trunc}}^2 \leq (\sigma_j^c)^2$
yields the stated bound.
\hfill$\square$
% \subsection{Proof of Corollary~\ref{cor:opt-displacement} (Optimization Displacement)}
% \label{app:proof-displacement}

% The displacement bound follows a three-term decomposition.

% \textbf{Step~(i): Quantile displacement.}
% Under Gaussianity of $X_j \mid c$, KLL rank error $\varepsilon_{\mathrm{sketch}}$
% implies that the empirical median $\hat{q}_{0.5,c}$ satisfies a probability-mass
% error of at most $\varepsilon_{\mathrm{sketch}}$ in rank. Inverting this through
% the class-conditional density, the spatial displacement obeys:
% \[
%   |\hat{q}_{0.5,c} - q_{0.5,c}|
%   \;\leq\; \frac{\varepsilon_{\mathrm{sketch}}}
%                 {p_{\min}(q_{0.5,c},\,\hat{q}_{0.5,c})},
% \]
% where $p_{\min}$ is the minimum Gaussian density over the interval between
% $q_{0.5,c}$ and $\hat{q}_{0.5,c}$.
% Under Gaussianity and bounded displacement, the density over this interval
% is bounded below by a value proportional to $1/(\sigma_c\sqrt{2\pi})$,
% giving $|\hat{q}_{0.5,c} - q_{0.5,c}| \leq \varepsilon_{\mathrm{sketch}}
% \cdot \sigma_c\sqrt{2\pi}$.

% \noindent\textit{Remark on bound direction.}
% Dividing by the \emph{maximum} density $1/(\sigma_c\sqrt{2\pi})$ at the
% Gaussian mode would yield a \emph{lower} bound on displacement, not an upper
% bound. The valid upper bound requires the \emph{minimum} density over the
% displaced interval; for small displacements under Gaussianity this equals
% $\sigma_c\sqrt{2\pi}$ asymptotically and the step above is valid in that
% regime. For large $|\zeta|$ where density is not approximately constant,
% this bound becomes loose. This is a standard asymptotic regime assumption
% in sketch-based quantile analysis~\cite{Karnin2016}; the practical
% impact is negligible as confirmed by the $\leq 4\times$ empirical safety
% margins in Table~\ref{tab:sketch-bound-kscale}.

% \textbf{Step~(ii): Threshold displacement.}
% Taking $\sigma_{\max} = \max_c \sigma_c$:
% \[
%   |\hat{v}^* - v^*| \leq \varepsilon_{\mathrm{sketch}} \cdot \sigma_{\max}\sqrt{2\pi}.
% \]

% \textbf{Step~(iii): Gini loss from displacement.}
% By Lemma~\ref{lem:gini-lipschitz} with Lipschitz constant
% $L_v = O(1/\sigma_{\min})$ and condition number
% $\kappa = \sigma_{\max}/\sigma_{\min}$:
% \[
%   F(S;\,v^*) - F(S;\,\hat{v}^*)
%   \leq L_v \cdot |\hat{v}^* - v^*|
%   \leq 4\kappa\,\varepsilon_{\mathrm{sketch}}.
% \]
% Adding the statistical term from Theorem~\ref{thm:composed} gives
% Eq.~\eqref{eq:composed-full}.\hfill$\square$

% \subsection{Proof of Proposition~\ref{prop:variance} (Post-Split Variance Reduction)}
% \label{app:proof-variance}

% \textit{Part~(i): Class probabilities.}
% The inherited prior is Dirichlet with total mass $S_0 = \alpha n_{\mathrm{parent}}$.
% After $t$ local samples, $S_t = S_0 + t$. The Dirichlet posterior variance:
% \[
%   \mathrm{Var}(\theta_c \mid \mathcal{D}_t)
%   = \frac{\beta_c^{(t)}(S_t - \beta_c^{(t)})}{S_t^2(S_t+1)}
%   \leq \frac{1}{4(S_t+1)} \leq \frac{1}{4(S_0+1)},
% \]
% where the first inequality uses the AM-GM bound
% $\beta_c^{(t)}(S_t - \beta_c^{(t)}) \leq S_t^2/4$.

% \textit{Part~(ii): Feature parameters.}
% Under the NIG conjugate model with $n_0 = \alpha n_c^{\mathrm{parent}}$
% pseudo-observations, the marginal posterior on $\mu_j^c$ at child creation
% is a Student-$t$ distribution with $\nu = n_0$ degrees of freedom and scale
% $(\sigma_j^c)^2/n_0$. Its variance exists when $n_0 > 2$ and equals:
% \[
%   \frac{\nu}{\nu-2}\cdot\frac{(\sigma_j^c)^2}{n_0}
%   = \frac{(\sigma_j^c)^2}{n_0-2}
%   = \frac{(\sigma_j^c)^2}{\alpha n_c^{\mathrm{parent}} - 2}.
% \]
% For the split feature, substituting
% $\sigma_{\mathrm{trunc}}^2 \leq (\sigma_j^c)^2$
% yields the stated bound. \hfill$\square$

\subsection{Full Derivation of Case~B (Cross-Routing)}
\label{app:case-b}

We provide the exact derivation for Case~B of Theorem~\ref{thm:tight},
where replacing sample $z_i$ with $z_i'$ changes the child to which it
is routed. The result is strictly tighter than the Case~A bound.

\begin{lemma}[Case~B Bound: Exact]
\label{lem:caseB}
Under Case~B (cross-routing): $z_i \in L$ has label~$1$,
replacement $z_i' \in R$ has label~$0$.
Then $|F(S) - F(S^{(i)})| \leq 2/n < 4/n$ exactly, for all finite $n$.
\end{lemma}

\begin{proof}
Let $m := n_{1,L}$ and $k := n_{1,R}$ be non-negative integers.
After the replacement: $L$ has $n_L-1$ samples with $m-1$ class-$1$,
and $R$ has $n_R+1$ samples with $k$ class-$1$.
Let $p = (m+k)/n$ be the class-$1$ proportion before the swap.

We compute $n\cdot\Delta F$ by expanding
$F = G_P - \frac{n_L}{n}G_L - \frac{n_R}{n}G_R$ into three terms.

\textit{Term~(1): parent Gini change.}
Using $G = 2p(1-p)$:
\[
  n\cdot\Delta G_P
  = n\!\left[2\!\left(p-\tfrac{1}{n}\right)\!\left(1-p+\tfrac{1}{n}\right) - 2p(1-p)\right]
  = 2(2p-1) - \frac{2}{n}.
\]

\textit{Term~(2): weighted left-child Gini change.}
Using $G_L = 2\frac{m}{n_L}\frac{n_L-m}{n_L}$:
\[
  n\cdot\Delta\!\left(\tfrac{n_L}{n}G_L\right)
  = 2(n_L-m)\!\left[\frac{m-1}{n_L-1} - \frac{m}{n_L}\right]
  = \frac{-2(n_L-m)^2}{n_L(n_L-1)}.
\]

\textit{Term~(3): weighted right-child Gini change.}
Using $G_R = 2\frac{k}{n_R}\frac{n_R-k}{n_R}$:
\[
  n\cdot\Delta\!\left(\tfrac{n_R}{n}G_R\right) = \frac{2k^2}{n_R(n_R+1)}.
\]

\textit{Exact formula:}
\begin{equation}
  n\cdot\Delta F
  = 2(2p-1) - \frac{2}{n}
    + \frac{2(n_L-m)^2}{n_L(n_L-1)}
    - \frac{2k^2}{n_R(n_R+1)}.
  \label{eq:caseB-exact}
\end{equation}

\textit{Bounding $|n\cdot\Delta F| < 2$.}
Write $T := \frac{n\cdot\Delta F}{2} = M(m) + K(k)$, where:
\begin{align}
  M(m) &:= \frac{2m-n_L-1}{n} + \frac{(n_L-m)^2}{n_L(n_L-1)}, \\
  K(k) &:= \frac{2k-n_R}{n} - \frac{k^2}{n_R(n_R+1)}.
\end{align}

$K(k)$ is a downward-opening quadratic in $k \in \{0,\ldots,n_R\}$
with maximum $K(k^*) = -n_R(n_L-1)/n^2 \leq 0$ and boundary values
satisfying $K(k) > -1$. Hence $-1 < K(k) \leq 0$.

$M(m)$ is an upward-opening quadratic in $m \in \{1,\ldots,n_L\}$
with minimum $M(m^*) = n_R(n_L-1)/n^2 \geq 0$ and boundary values
satisfying $M(m) < 1$. Hence $0 \leq M(m) < 1$.

Therefore $-1 < K(k) + M(m) < 1$, giving
$|T| < 1$ and $|n\cdot\Delta F| < 2 < 4$.
\end{proof}

\begin{remark}[Extension to $K > 2$]\label{rem:caseb-multiclass}
In Case~B with $K > 2$, the replacement $z_i \to z_i'$ moves one
sample from child $L$ to child $R$ (or vice versa), changing exactly
two class counts: $n_{y_i, L}$ decreases by 1 and $n_{y_i', R}$
increases by 1, while $n_L$ and $n_R$ change by $\mp 1$.
The Gini split gain depends on child counts only through
$\sum_c p_{c,s}^2 = \sum_c (n_{c,s}/n_s)^2$.
Perturbing one class count $n_{c_0, L}$ by $-1$ and $n_L$ by $-1$:
\[
  \Delta\!\left(\sum_c \frac{n_{c,L}^2}{n_L^2}\right)
  = \frac{(n_{c_0,L}-1)^2 + \sum_{c \neq c_0} n_{c,L}^2}{(n_L-1)^2}
  - \frac{\sum_c n_{c,L}^2}{n_L^2}.
\]
Direct computation shows $|n \cdot \Delta F| < 2$ holds for all
non-negative integer counts by the same bounding argument as the
$K=2$ case: the numerator change is bounded by $2n_L - 1$ and the
combined effect satisfies $|T| = |M(m) + K(k)| < 1$.
The binary case ($K=2$) is thus the binding extremal configuration,
confirming that $4/n$ is tight for all $K$.
\end{remark}
% \begin{remark}
% The bound $|n\cdot\Delta F| \leq 2 < 4$ is strictly tighter than Case~A.
% The global tightness witness of Theorem~\ref{thm:tight} is therefore a
% Case~A configuration, confirming that $4/n$ is the tight constant for the
% overall split gain sensitivity.
% For $K > 2$, cross-routing moves sample $z_i$ from child $L$ to child $R$
% (or vice versa), altering the counts of at most two classes simultaneously
% (the label of $z_i$ and the label of $z_i'$). The Gini gain $\Delta F$
% depends on $n_{c,L}$ only through the quadratic $\sum_c p_{c,L}^2$;
% the same $M(m) + K(k)$ decomposition above applies class-by-class, so
% $|n\cdot\Delta F| < 2$ holds for arbitrary $K$.
% \end{remark}

%%%%%%%%%%%%%%%%%%%%%%%%%%%%%%%%%%%%%%%%%%%%%%%%%%%%%%%%%%%%

% \newpage
% \input{checklist.tex}


\end{document}
